a% !TEX root = ../EDRmain.tex
%\def\manifold{\operatorname{ma}}
\def\manifold{\mathcal{M}}
\def\subM{{_{\! \manifold}}}

\def\numb{N}
\def\numbd{\numb^{\circ}}
\def\nmin{\numb_{\operatorname{loc}}}
\def\nman{\numb_{_{\! \manifold}}}

\def\Kern{\mathcal{K}}
\def\sigmaX{\sigma_{\!_{X}}}
\def\Xbar{\Mv}
\def\corr{\tau}

\def\LL{\mathcal{L}}

\def\jd{j^{\circ}}
\def\jjl{j\jl}
\def\jl{\alpha}
\def\jld{\jl^{\circ}}

%\def\weightM{\omega}
\def\weightM{\delta}
\def\WeightM{\Delta}

\def\svd{\sv^{\circ}}
\def\xvm{\xv^{\dag}}

\def\JJJ{\mathcal{J}}
\def\nJJJ{\numb_{\!_{\JJJ}}}
\def\JJJd{\JJJ^{\circ}}
\def\init{\operatorname{in}}
\def\outp{\operatorname{out}}
\def\IEDR{\mathcal{A}}


%%%%%%%%%%%%%%%%%%%%%%%%%%%%%%%  new section  %%%%%%%%%%%%%%%%%%%%%%%%%%%%%%%
\Section{A single-index model}
\label{Ssingind}

Consider the e.d.r. regression problem
\begin{EQA}
	Y_{i}
	&=&
	\fs(\betav^{\T} \Xv_{i}) + \eps_{i} \, ,
\label{YifXieiedrSI}
\end{EQA}
where \( \Xv_{i} \in \R^{\dimd} \) for large \( \dimd \), \( \betav \) is an index vector and
\( \fs \) is a univariate \emph{link function}.
Introduce also a function \( \gs(x) = \E \bigl( Y \cond \Xv=\xv \bigr) \).
A structural single-index assumption yields the relations
\begin{EQA}
	\gs(\xv)
	&=&
	\fs(\betav^{\T} \xv) ,
	\\
	\nabla \gs(\xv)
	&=&
	\fs'(\betav^{\T} \xv) \betav.
\label{fxgKxedrSI}
\end{EQA} 
In particular, the derivative of \( \gs \) at any point \( \xv \) is proportional
to the target vector \( \betav \).
This suggests recovering \( \betav \) from the nonparametrically estimated gradient of \( \gs \).


%%%%%%%%%%%%%%%%%%%% section %%%%%%%%%%%%%%%%%%%%
\Subsection{Preliminaries}
%Define \( \Psi(\xv) = \bigl( \psi_{m}(\xv) \bigr) \in \R^{\dimd+1} \) for \( \psi_{0}(\xv) \equiv 1 \), 
%\(	\psi_{m}(\xv) =	\langle \xv ,\ev_{m} \rangle \), \( m = 1,\ldots,\dimd \). 
For a collection of centers \( \{ \xv_{j} \}_{j \in \JJJ} \) with \( |\JJJ| = \nJJJ \),
localizing tensor \( \TEDR \), 
and a localizing kernel \( \Kern \), consider localizing weights
\begin{EQA}
	\weight_{ij}
	&=&
	\Kern\bigl( \| \TEDR(\Xv_{i} - \xv_{j}) \|^{2} \bigr) 
	\, ,
	\qquad
	i=1,\ldots,n,
	\quad
	j \in \JJJ
	\, .
\label{s7cuqkdwywdhbjwbfjwbdj}
\end{EQA}
At the first step, we apply \( \TEDR = \hiso^{-2} \Id_{\dimp} \)
leading to isotropic weights 
\( \weight_{ij} = \Kern\bigl( \hiso^{-2} \| \Xv_{i} - \xv_{j} \|^{2} \bigr) \).
The bandwidth \( \hiso \) has to be selected to ensure sufficiently many points in almost every vicinity \( U_{\hiso}(\xv_{j}) \) of \( \xv_{j} \)
of radius \( \hiso \). 

Further, for each \( j \in \JJJ \), let \( \Sph_{j} \) be a collection of 
directions (unit vectors) \( \sph \in \R^{\dimd} \) with \( |\Sph_{j}| = \nSph \).
Define for each \( j \in \JJJ \) and \( \sph \in \Sph_{j} \)
\begin{EQ}[rcl]
	\Ima_{j,\sph}
	&=&
	\sumi Y_{i} \langle \Xv_{i} - \xv_{j} \, ,\sph \rangle \, \weight_{ij} 
	\, ,
	\qquad
	S_{j,\sph}
	=
	\sumi \langle \Xv_{i} - \xv_{j} \, ,\sph \rangle \, \weight_{ij}
	\, ,
%\label{ytd7e7ydwuhd5eythdY}
%\end{EQ}
%and also the vector \( \Uv_{j,\sph} \) in \( \R^{\dimd} \)
%\begin{EQA}[rcl]
	\\
	\Uv_{j,\sph}
	&=&
	\sumi (\Xv_{i} - \xv_{j}) \, \langle \Xv_{i} - \xv_{j} \, ,\sph \rangle \, \weight_{ij}
	\in \R^{\dimd}
	\, .
	\qquad
\label{ytd7e7ydwuhd5eythdm}
\end{EQ}
%where
%\begin{EQA}[c]
%	N_{j}
%	\eqdef
%	\sumi \weight_{ij}
%	\, ,
%	\quad
%	\Sv_{j}
%	\eqdef
%	\sumi (\Xv_{i} - \xv_{j}) \weight_{ij}
%	\, ,
%	\quad
%	\VP_{j}
%	\eqdef
%	\sumi (\Xv_{i} - \xv_{j})(\Xv_{i} - \xv_{j})^{\T} \weight_{ij}
%	\, .
%\end{EQA}
Model equation \eqref{YifXieiedrSI} implies
\begin{EQA}
	\Ima_{j,\sph}
	&=&
	\sumi Y_{i} \langle \Xv_{i} - \xv_{j} \, ,\sph \rangle \, \weight_{ij}  
	\\
	&=&
	\sumi \fs(\betav^{\T} \Xv_{i}) \langle \Xv_{i} - \xv_{j} \, ,\sph \rangle \, \weight_{ij}
	+ \sumi \eps_{i} \langle \Xv_{i} - \xv_{j} \, ,\sph \rangle \, \weight_{ij} \, .
\label{hgghjkklvddfrdshwey}
\end{EQA}
Further, localization by the weights 
\( \weight_{ij} \) enables us to locally approximate \( \fs(t) \) 
in a vicinity of \( t_{j} = \langle \betav, \xv_{j} \rangle \) by a linear function
\( \fcn_{j} + \fcl_{j} (t - t_{j}) \).
This leads to the approximation 
\begin{EQA}
	&& \nquad
	\sumi  \fs(\betav^{\T} \Xv_{i}) \langle \Xv_{i} - \xv_{j} \, ,\sph \rangle \, \weight_{ij} 
	\approx 
	\sumi \bigl\{ \fcn_{j} + \fcl_{j} (\Xv_{i} - \xv_{j})^{\T} \betav \bigr\} \langle \Xv_{i} - \xv_{j} \, ,\sph \rangle \, \weight_{ij} 
	\\
	&=&
	\fcn_{j} S_{j,\sph} + \fcl_{j} \langle \Uv_{j,\sph} \, ,\betav \rangle
%	\KSe_{j} \binom{\fcn_{j}}{\fcl_{j} \, \betav} 
%	=
%	\begin{pmatrix}
%		\fcn_{j} N_{j} + \fcl_{j} \Sv_{j}^{\T} \betav
%		\\
%		\fcn_{j} \Sv_{j} + \fcl_{j} \VP_{j} \betav
%	\end{pmatrix}
	\, .
\label{yd6x5ewbduwuqf64wfgfy6}
\end{EQA}
The calming device suggests introducing additional parameters (observables) 
\( \imav_{j,\sph} = \E \Ima_{j,\sph} \).
The full dimensional parameter vector \( \prmtv \) includes
the target parameter \( \betav \), the observables \( \imav = (\imav_{j,\sph}) \in \R^{\nJJJ \times \nSph} \), 
and the coefficients \( \thetav = (\fcn_{j},\fcl_{j})_{j \in \JJJ} \in \R^{2\nJJJ} \).
The corresponding objective function reads
\begin{EQA}
	\LL(\prmtv)
	&=&
	\LL(\betav,\thetav,\imav)
	=
	\frac{1}{2} \| \Imav - \imav \|^{2} 
	+ \frac{1}{2} \sum_{j \in \JJJ} \sum_{\sph \in \Sph_{j}} 
		\bigl( \ima_{j,\sph} - \fcn_{j} S_{j,\sph} - \fcl_{j} \langle \Uv_{j,\sph} \, ,\betav \rangle \bigr)^{2} 
	\\
	&=&
	\frac{1}{2} \sum_{j \in \JJJ} 
	\bigl( \| \Imav_{j} - \imav_{j} \|^{2} + \bigl\| \imav_{j} - \fcn_{j} \Sv_{j} - \fcl_{j} \, \UV_{j} \betav \bigr\|^{2} \bigr) 
	\, ,
\label{gtgyuy6ew5dtde32322jt}
\end{EQA}
where \( \imav_{j} = (\ima_{j,\sph} \, , \sph \in \Sph_{j}) \in \R^{\nSph} \),
\( \Sv_{j} = (S_{j,\sph} \, , \sph \in \Sph_{j}) \in \R^{\nSph} \),
and \( \UV_{j} \) is a \( \nSph \times \dimd \)-matrix with rows \( \Uv_{j,\sph}^{\T} \), that is,
\( \UV_{j}^{\T} = (\Uv_{j,\sph})_{\sph \in \Sph_{j}} \).
Due to the product \( \fcl_{j} \, \UV_{j} \betav \), this objective function is not quadratic in the scope of parameters.
More precisely, it is a polynomial of order 4, and it can be considered as a special case of the error-in-operator model.
The challenge for a careful analysis is a lack of identifiability: rescaling the vector \( \betav \) by any constant
\( c \) and all the coefficients \( \fcl_{j} \) by \( c^{-1} \) does not change the model equation. 
To overcome this problem and enforce identifiability, we require \( \| \betav \| = 1 \).

\paragraph{Elimination of the observables \( \zv \).}
Representation \eqref{gtgyuy6ew5dtde32322jt} is useful for theoretical analysis.
Numerically, one can eliminate the auxiliary variables \( \zv = (\zv_{j,\sph}) \).

\begin{lemma}
For any \( j \in \JJJ \) and any
\( \fcn_{j},\fcl_{j} \, , \betav \), it holds
\begin{EQA}[c]
	\inf_{\imav_{j} \in \R^{\nSph}} \Bigl( 
		\| \Imav_{j} - \imav_{j} \|^{2} 
		+ \bigl\| \imav_{j} - \fcn_{j} \Sv_{j} - \fcl_{j} \, \UV_{j} \betav \bigr\|^{2} 
	\Bigr) 
	=
	\frac{1}{2} \bigl\| \Imav_{j} - \fcn_{j} \Sv_{j} - \fcl_{j} \, \UV_{j} \betav \bigr\|^{2}
	\, .
\end{EQA}
\end{lemma}
{\color{blue} Please, check.}


This result reduces problem \eqref{gtgyuy6ew5dtde32322jt} to
\begin{EQA}[c]
	\sum_{j \in \JJJ} \bigl\| \Imav_{j} - \fcn_{j} \Sv_{j} - \fcl_{j} \, \UV_{j} \betav \bigr\|^{2}
%	+ \lambda \| \betav - \betav_{\init} \|^{2}
	\to
	\inf_{\betav,(\fcn_{j},\fcl_{j})_{j \in \JJJ}}
	\, ,
\label{ikdfjoiefjoieuvbjmr3e}
\end{EQA}
where \( \betav \) is a point on the unit sphere \( \Sphere_{\dimd} \) in \( \R^{\dimd} \).

\paragraph{Relaxation and regularization}
The set \( \Sphere_{\dimd} \) is not convex in \( \R^{\dimd} \), 
and minimization over this set is numerically hard.
We apply the usual relaxation trick by allowing at every step of the minimization procedure 
an unconstrained optimization over \( \betav \in \R^{\dimd} \) and then by projecting 
the obtained solution to the unit sphere.
Additional numerical stability can be ensured 
by a penalty \( \mfrac{\lambda}{2} \| \betav - \betav_{\init} \|^{2} \) with 
\( \betav_{\init} \in \R^{\dimd} \)
meaning a starting guess or the previous step estimate.
%One can use \( \betav_{\init} = (1,0,\ldots,0)^{\T} \) or a random unit vector un \( \R^{\dimd} \).
%It is only important that \( \langle \betav,\ev_{1} \rangle \) significantly deviates from zero.
With a tuning parameter \( \lambda \), this leads to minimization of
\begin{EQA}
	\LL_{\lambda}(\prmtv)
	&=&
	\frac{1}{2} \sum_{j \in \JJJ} \bigl\| \Imav_{j} - \fcn_{j} \Sv_{j} - \fcl_{j} \, \UV_{j} \betav \bigr\|^{2} 
	+ \frac{\lambda}{2} \| \betav - \betav_{\init} \|^{2} 
	\, 
\label{tjjc7c6f5fd5e5erjwyhfd}
\end{EQA}
at one step of the procedure.
%

%%%%%%%%%%%%%%%%%%%% section %%%%%%%%%%%%%%%%%%%%
\Subsection{Alternating optimization}
\label{SLLaltern}
Now we describe the alternating optimization procedure for solving \eqref{ikdfjoiefjoieuvbjmr3e}
with given \( \Imav_{j} \), \( \Sv_{j} \), and \( \UV_{j} \), \( j \in \JJJ \).
The idea is to iterate optimization of the objective function w.r.t. 
\( (\fcn_{j},\fcl_{j})_{j \in \JJJ} \) for \( \betav \) fixed,
and w.r.t. \( \betav \) for \( (\fcn_{j},\fcl_{j})_{j \in \JJJ} \) fixed.
Note that for a fixed \( \betav \), the optimization problem \eqref{ikdfjoiefjoieuvbjmr3e}
can be decomposed into small subproblems 
\begin{EQA}[c]
%	\inf_{(\fcn_{j},\fcl_{j})} 
	\bigl\| \Imav_{j} - \fcn_{j} \Sv_{j} - \fcl_{j} \, \UV_{j} \betav \bigr\|^{2}
	\to \inf_{\fcn_{j},\fcl_{j}}
	\, ,
	\qquad
	j \in \JJJ
	\, .
\end{EQA}
Each of them is quadratic in \( \fcn_{j},\fcl_{j} \) and can be solved analytically.

\begin{lemma}
\label{Lcjlj}
It holds 
\begin{EQA}[c]
	\arginf_{\fcn_{j},\fcl_{j}} 
	\bigl\| \Imav_{j} - \fcn_{j} \Sv_{j} - \fcl_{j} \, \UV_{j} \betav \bigr\|^{2}
	=
	\begin{pmatrix}
		\| \Sv_{j} \|^{2} 
		& 
		\langle \Sv_{j} \, , \UV_{j} \betav \rangle
		\\
		\langle \Sv_{j} \, , \UV_{j} \betav \rangle  
		& 
		\| \UV_{j} \betav \|^{2} 
	\end{pmatrix}^{-1} 
	\binom{ \langle \Imav_{j} \, , \Sv_{j} \rangle }
	{\langle \Imav_{j} \, , \UV_{j} \betav \rangle }
	\, .
	\qquad
\label{ihi88eufiut545yhfee}
\end{EQA}
\end{lemma}
{\color{blue} please check.}

Similarly, for fixed \( \{ (\fcn_{j},\fcl_{j}), j \in \JJJ \} \), minimization of \eqref{ikdfjoiefjoieuvbjmr3e} w.r.t. \( \betav \)
is quadratic and can be solved analytically.

\begin{lemma}
\label{Lbeta}
It holds 
\begin{EQA}[rcl]
	&& \nquad
	\arginf_{\betav \in \R^{\dimd}} \biggl\{ \sum_{j \in \JJJ} \bigl\| \Imav_{j} - \fcn_{j} \Sv_{j} - \fcl_{j} \, \UV_{j} \betav \bigr\|^{2}
	+ \lambda \| \betav - \betav_{\init} \|^{2} 
	\biggr\}	
	\\
	&=&
	\biggl\{ \sum_{j \in \JJJ} \fcl_{j}^{2} \, \UV_{j}^{\T} \UV_{j} + \lambda \Id_{\dimd} \biggr\}^{-1}
	\sum_{j \in \JJJ} \bigl\{ \fcl_{j} \, \UV_{j}^{\T} (\Imav_{j} - \fcn_{j} \Sv_{j}) 
	+ \lambda \betav_{\init}
	\bigr\}
	\, .
\label{uh8888iiftbe7fuhfff}
\end{EQA}
{\color{blue} please check}
\end{lemma}


\paragraph{One-step alternating}
The procedure assumes a starting guess \( \betav_{\init} \) to be given.
%
%\paragraph{Initialization.}
%Set \( \betav_{0} = \betav_{\init} \) and \( m=1 \).
%
%\paragraph{Step \( m \geq 1 \). Assumed \( \betav_{\init} \) available.}
\begin{mylist}
\item 
Compute for every \( j \in \JJJ \)
\begin{EQA}[c]
	(\hat{\fcn}_{j},\hat{\fcl}_{j})
	=
	\arginf_{\fcn_{j},\fcl_{j}} \bigl\| \Imav_{j} - \fcn_{j} \Sv_{j} - \fcl_{j} \, \UV_{j} \betav_{\init} \bigr\|^{2}
	\, ;
\end{EQA}
see \eqref{ihi88eufiut545yhfee} of Lemma~\ref{Lcjlj}.

\item 
Then with \( (\hat{\fcn}_{j},\hat{\fcl}_{j})_{j \in \JJJ} \), update \( \betav \)
by \eqref{uh8888iiftbe7fuhfff} of Lemma~\ref{Lbeta}:
\begin{EQA}[c]
	\hat{\betav}
	=
	\arginf_{\betav} \biggl\{ \sum_{j \in \JJJ} \bigl\| \Imav_{j} - \hat{\fcn}_{j} \Sv_{j} - \hat{\fcl}_{j} \, \UV_{j} \betav \bigr\|^{2}
	+ \lambda \| \betav - \betav_{\init} \|^{2} 
	\biggr\}	
	\, .
\end{EQA}


\item
Normalize \( \hat{\betav} \) to one: 
\begin{EQA}[c]
	\betav_{\outp} 
	= 
	\mfrac{1}{\| \hat{\betav} \|} \, \hat{\betav} 
	\, .
\end{EQA}
%With \( \rr_{m} \eqdef \| \betav_{m} \| \), update \( \betav_{m} \to \betav_{m}/\rr_{m} \) and \( \fcl_{j} \to \rr_{m} \fcl_{j} \) for
%all \( j \in \JJJ \). 

\end{mylist}

\paragraph{Alternating optimization}
One-step alternating can be repeated several times.
At every iteration, the output \( \betav_{\outp} \) of the previous step is used as input \( \betav_{\init} \).
The procedure delivers the last step unit vector \( \betav_{\outp} \) after \( k_{\max} \) iterations.


{\color{blue} check one-step improvement and accuracy of \( \betav_{\outp} \) after \( k \) iterations.}



%%%%%%%%%%%%%%%%%%%% section %%%%%%%%%%%%%%%%%%%%
\Subsection{Anisotropic weights and structural adaptation}
\label{SanisoSI}
The local weights \( \weight_{ij} \) from \eqref{s7cuqkdwywdhbjwbfjwbdj} involve a local tensor \( \TEDR \).
The single-index model structure \( \gs(\xv) = \fs(\betav^{\T} \xv) \) in \eqref{YifXieiedrSI} suggests
anisotropic weights
\begin{EQA}
	\weights_{ij}
	&=&
	\Kern\bigl( h^{-2} |(\Xv_{i} - \xv_{j})^{\T} \betavs|^{2} \bigr) .
\label{s7cuqkdwywdhbjwbfjwbdji}
\end{EQA}
Unfortunately, such defined weights require knowledge of the index vector \( \betavs \) and 
a correctly specified single-index model. 
We follow the structure-adaptive approach from \cite{HJPS}, which applies anisotropic weights 
with an increasing anisotropy parameter during iterations.
As mentioned earlier, for initialization, we apply a non-informative choice 
\begin{EQA}[c]
	\weight_{ij} = \Kern\bigl( \hiso_{0}^{-2} \| \Xv_{i} - \xv_{j} \|^{2} \bigr) 
	\, .
\end{EQA}
The bandwidth \( \hiso_{0} \) can be selected as the smallest value ensuring
\begin{EQA}[c]
	\frac{1}{\nJJJ}\sum_{j \in \JJJ} \sumi \Kern\Bigl( \mfrac{\| \Xv_{i} - \xv_{j} \|^{2}}{\hiso_{0}^{2}} \Bigr)
	\geq 
	\nmin
	\, .
\label{jhgiverkytde4w34etycj}
\end{EQA}
This condition means that in a local vicinity of each point \( \xv_{j} \) there are (on average) 
at least \( \nmin \) design points \( \Xv_{i} \).

Anisotropic weights assume some prior information about the structure of the model, which is accumulated 
in the localizing tensor \( \TEDR \).
Namely, for two values \( \hiso \) and \( \aniso \leq 1 \) and a current guess \( \betav \), anisotropic weights \( \weight_{ij} \)
correspond to the tensor 
\begin{EQA}[rcl]
	\TEDR^{2}
	&=&
	\hiso^{-2} \bigl( \aniso^{2} \, \Id_{\dimd} + \betav \betav^{\T} \bigr)
\label{iufhvu8kjbmnll475edf}
\end{EQA}
leading to 
\begin{EQA}[c]
	\weight_{ij}
	=
	\Kern\bigl( \| \TEDR (\Xv_{i} - \xv_{j}) \|^{2} \bigr)
	=
	\Kern\bigl( 
		\mfrac{\aniso^{2} \| \Xv_{i} - \xv_{j} \|^{2}}{\hiso^{2}} 
		+ \mfrac{\langle \Xv_{i} - \xv_{j}, \betav \rangle^{2}}{\hiso^{2}} 
	\bigr)
	\, .
\end{EQA}
The quantity \( \aniso \) describes the anisotropy strength.
Small values of \( \hiso, \aniso \) correspond to strong localization in direction \( \betav \)
and mild localization in orthogonal directions. 
During iteration, the procedure decreases \( \hiso \) and \( \aniso \) in a way 
to ensure sufficiently many design points in the local vicinity of \( \xv_{j} \) 
for almost every \( j \in \JJJ \).
For fixed \( \hiso \) and \( \betav \), the anisotropy parameter \( \aniso \) is the smallest value ensuring
\begin{EQA}[c]
	\frac{1}{\nJJJ} \sum_{j \in \JJJ} 
	\sumi \Kern\bigl( \| \TEDR (\Xv_{i} - \xv_{j}) \|^{2} \bigr)
%	=
%	\frac{1}{\nJJJ} \sum_{j \in \JJJ} \sumi \Kern\Bigl( 
%		\mfrac{\aniso^{2} \| \Xv_{i} - \xv_{j} \|^{2}}{\hiso^{2}} + \mfrac{\langle \Xv_{i} - \xv_{j}, \betav \rangle^{2}}{\hiso^{2}} 
%	\Bigr)
%	=
%	\frac{1}{\nJJJ} \sum_{j \in \JJJ} \sumi \weight_{ij}
	\geq 
	\nmin
	\, .
	\qquad
\label{jhgiverkytde4w34etycja}
\end{EQA}
It is easy to see that a decay of \( \hiso \) leads to a decay of \( \aniso \).

%%%%%%%%%%%%%%%%%%%% section %%%%%%%%%%%%%%%%%%%%
\Subsection{Initialization by local linear estimation}
\label{Siniloclin}
The objective function \( \LL(\prmtv) \) is not convex. 
Therefore, the starting direction \( \betav_{\init} \) is important for a sensitive final result.
\cite{HJPS} suggested fixing \( \betav_{\init} \)
by local linear estimation of the gradient of \( \gs \) at the points \( \xv_{j} \).
Namely, fix \( \hiso_{\lin} \) as the smallest value ensuring 
\begin{EQA}[c]
	\frac{1}{\nJJJ}\sum_{j \in \JJJ} 
	\sumi \Kern\Bigl( \mfrac{\| \Xv_{i} - \xv_{j} \|^{2}}{\hiso_{\lin}^{2}} \Bigr)
	\geq 
	\numb_{\lin}
	\, .
\label{hhfiecredgvhihiue86t}
\end{EQA}
%The default value of \( \numb_{\lin} \) is \( n/\nmin \).
The proposal requires \( \numb_{\lin} > \dimd + 1 \).
Now, for each \( j \in \JJJ \), compute
\begin{EQA}[c]
	\begin{pmatrix}
		\hat{\gs}(\xv_{j}) \\
		\hat{\nabla \gs}(\xv_{j})
	\end{pmatrix}
	=
	\biggl\{ \sumi \binom{1}{\Xv_{i} - \xv_{j}} \, \binom{1}{\Xv_{i} - \xv_{j}}^{\T}
		\Kern\Bigl( \mfrac{\| \Xv_{i} - \xv_{j} \|^{2}}{\hiso_{\lin}^{2}} \Bigr)
	\biggr\}^{-1}
	\sumi Y_{i} \, \binom{1}{\Xv_{i} - \xv_{j}} \, 
		\Kern\Bigl( \mfrac{\| \Xv_{i} - \xv_{j} \|^{2}}{\hiso_{\lin}^{2}} \Bigr)
	\, .
\end{EQA}
Define \( \betav_{\init} \) as the first principal direction
for the set of the estimated gradients \( \{ \hat{\nabla \gs}(\xv_{j}) \}_{j \in \JJJ} \).
For this, compute
\begin{EQA}[c]
	\JEDR
	\eqdef
	\sum_{j \in \JJJ} \hat{\nabla \gs}(\xv_{j}) \{ \hat{\nabla \gs}(\xv_{j}) \}^{\T}
	\, .
\end{EQA}
Then \( \betav_{\init} \) is the the eigenvector of \( \JEDR \) corresponding to its largest eigenvalue
\( \lambda_{1}(\JEDR) \).
If \( \lambda_{1}(\JEDR) > \lambda_{2}(\JEDR) \) then this eigenvector is uniquely defined up to direction sign.
However, we only need the product \( \betav_{\init} \, \betav_{\init}^{\T} \) which is uniquely defined.

\begin{mylist}
	\item
	{\color{blue} check significance \( \cos\langle \betavs,\betav_{\init} \rangle \gg 0 \).}
	\item
	{\color{blue} explore the impact of \( \numb_{\lin} \) from \eqref{hhfiecredgvhihiue86t}
	in the range between \( 2\dimd \) and \( n/\nmin \).}
	\item
	{\color{blue} explore the impact of \( \nJJJ \).}
\end{mylist}

%%%%%%%%%%%%%%%%%%%% section %%%%%%%%%%%%%%%%%%%%
\Subsection{Design of experiments}
Fix the values \( d \) and \( n \) such that \( n \geq 10 d \).
Also, fix the values \( \sigma_{\eps} \), \( \sigmaX \).
%\( \nJJJ \leq n \), \( \nSph \leq \dimd \), and \( \corr \in [0,1) \).
Finally, fix a unit vector \( \betavs \) and a univariate function \( \fs(x) \); e.g \( f(x) = x^{2} \).
\begin{mylist}

\item 
generate a correlated design \( \Xv_{i} = \sigmaX \bigl\{ \corr \gaussv_{0} + (1 - \corr) \gaussv_{i} \bigr\} \),
where \( \gaussv_{i} \) are i.i.d. standard Gaussian in \( \R^{\dimd} \), \( i=0,1,\ldots,n \);
\item 
generate errors \( \eps_{i} \sim \ND(0,\sigma_{\eps}^{2}) \) and compute \( Y_{i} = \fs(\Xv_{i}^{\T} \betavs) + \eps_{i} \);

%\item
%generate \( \betav_{\ini} \) uniformly on the unit sphere in \( \R^{\dimd} \).
\end{mylist}

%%%%%%%%%%%%%%%%%%%% section %%%%%%%%%%%%%%%%%%%%
\Subsection{Tuning parameters}
For the procedure, we have to fix a set of tuning parameters; see Table~\ref{TdefSI}
for the complete list and default values.
%
\begin{table}[h]
\begin{tabular}{c|c|c}
        \hline
        Parameter & Meaning & default \\
        \hline
        \( \nmin \) & \( \# \) of design points \( \Xv_{i} \) in a local vicinity & 10 \\
        \( \numb_{\lin} \) & parameter for local linear estimation & 
        between \( 2d \) and \( n/\nmin \) \\
        \( \nJJJ \) & size of \( \JJJ \) 				& between \( n/\nmin \) and \( n \) \\
		\( \nSph \) & size of each \( \Sph_{j} \)		& \( \min\{ \nmin, \dimd \} \) \\
%\( a (= 2^{1/\dimm}) \), 
%		\( \lambda & (= \nmin \dimm) \),
		\( \lambda \) & Lagrange multiplier 		& between 1 and \( \nmin \) \\
		\( \Kern \) & the kernel & Epanechnikov: \( \Kern(t) = (1-t^{2})_{+} \) \\
		\( k_{\max} \) & the number of repetitions in AO & 5 \\
		\( a \)	& decrease factor for \( \hiso_{k} \) & \( \sqrt{2} \) \\
		\( \hiso_{\min} \) & the smallest \( \hiso_{k} \) & \( 10 \sigmaX/n \) \\
		\hline
    \end{tabular}
\caption{Tuning parameters for single-index procedure}    
\label{TdefSI}
\end{table}
%
%\begin{mylist}
%
%\item 
%the value \( \nmin \gg 1 \); by default, \( \nmin = 10 \);
%
%\item 
%the value \( \nJJJ \gg 1 \); by default, \( \nJJJ = n \);
%
%\item
%the value \( \nSph \gg 1 \); by default, \( \nSph = \min\{ \nmin, \dimd \} \);
%
%\item 
%the kernel \( \Kern \); by default, Epanechnikov kernel \( \Kern(t) = (1-t^{2})_{+} \);
%
%\item 
%the Lagrange penalizing constant \( \lambda \), by default, 
%\( \lambda = 1 \);
%
%\item 
%the factor \( a > 1 \); by default, \( a = \sqrt{2} \);
%
%\item 
%the value \( \hiso_{\min} > 1 \); by default, \( \hiso_{\min} = 10 \sigmaX/n \).
%\end{mylist}
%
Parameters \( \nJJJ \) and \( \nSph \) are mainly for controlling the numerical
complexity of the procedure.
Parameter \( \lambda \) has the meaning of the step size, 
it controls stability and convergence of alternating optimization
from Section~\ref{SLLaltern}.
Parameter \( \nmin \) is essential for the structure-adaptive block.
However, the procedure exhibits a very stable performance 
for a minor variation of the parameters around the default values.

%%%%%%%%%%%%%%%%%%%% section %%%%%%%%%%%%%%%%%%%%
\Subsection{Procedure}

\paragraph{Initialization.}
\begin{mylist}

\item 
Fix a collection of points \( (\xv_{j})_{j \in \JJJ} \).
As an example, for each \( \Xv_{i} \), 
decide by Bernoulli experiment with success probability \( \theta \eqdef \nJJJ/n \), 
whether this point is used for 
constructing the center \( \xv_{j} \) by \( \xv_{j} = \Xv_{i} + \sigmaX \gaussv_{j} \);
ensure \( \#(\xv_{j}) = \nJJJ \);
%\item
%for each \( j \in \JJJ \), generate the set \( \Sph_{j} \) of unit vectors \( \sph \in \R^{\dimd} \)
%with \( |\Sph_{j}| = \nSph \).
%One can use random draws \( \sph \eqd \gaussv/\| \gaussv \| \) for \( \gaussv \sim \ND(0,\Id_{\dimd}) \);

\item
Set \( k=0 \).

\item
Fix \( \betav_{\init} \) using local linear regression as
in Section~\ref{Siniloclin}.

\item
Fix \( \hiso_{0} \) as the smallest value ensuring \eqref{jhgiverkytde4w34etycj}.
Define \( \TEDR_{0} = \hiso_{0}^{-1} \Id_{\dimp} \).
%\begin{EQA}[c]
%	\frac{1}{\nJJJ} \sum_{j \in \JJJ} 
%	\sumi \Kern\Bigl( \mfrac{\| \Xv_{i} - \xv_{j} \|^{2}}{\hiso_{0}^{2}} \Bigr)
%%	\bigl( \hiso_{0}^{-2} \| \Xv_{i} - \xv_{j} \|^{2} \bigr)
%	\geq 
%	\nmin
%	\, .
%\end{EQA}
%Also, set \( \haniso_{0} = \hiso_{0} \).

%\item 
%For each \( j \in \JJJ \) and \( \sph \in \Sph_{j} \), 
%compute \( \Ima_{j,\sph}^{(0)} \, , S_{j,\sph}^{(0)} \, , \Uv_{j,\sph}^{(0)} \) by \eqref{ytd7e7ydwuhd5eythdm}
%with 
%\begin{EQA}[c]
%	\weight_{ij} = \Kern\bigl( \hiso_{0}^{-2} \| \Xv_{i} - \xv_{j} \|^{2} \bigr) 
%	\, .
%\end{EQA}
%\item 
%%Find \( \prmtv^{(0)} = \{ \betav_{\init},(\fcn_{j}^{(0)},\fcl_{j}^{(0)}) \} \) as minimizer of 
%%\( \LL_{\lambda}(\prmtv) \) from \eqref{tjjc7c6f5fd5e5erjwyhfd} 
%With 
%\( \Imav_{j}^{(0)} = \bigl( \Ima_{j,\sph}^{(0)} \bigr)_{\sph \in \Sph_{j}} \),
%\( \Sv_{j}^{(0)} = \bigl( S_{j,\sph}^{(0)} \bigr)_{\sph \in \Sph_{j}} \),
%\( \UV_{j}^{(0)} = (\Uv_{j,\sph}^{(0)} \, , \sph \in \Sph_{j})^{\T} \),
%solve
%\begin{EQA}[c]
%	\bigl\{ \betav_{\init},(\fcn_{j}^{(0)},\fcl_{j}^{(0)})_{j \in \JJJ} \bigr\}
%	=
%	\arginf_{\betav \, , (\fcn_{j},\fclv_{j})_{j \in \JJJ}}
%	\sum_{j \in \JJJ} \Bigl( \| \Imav_{j}^{(0)} - \fcn_{j} \Sv_{j}^{(0)} - \fcl_{j} \, \UV_{j}^{(0)} \betav \bigr\|^{2} \Bigr) 
%%	+ \frac{\lambda}{2} \| \betav - \betav_{\ini} \|^{2}
%	\, 
%	\qquad
%\end{EQA}
%by alternating optimization with a random start \( \betav_{\init} \); see Section~\ref{SLLaltern}.

%\item 
%%:
%Rescale \( \betav_{\init} \) by \( \| \betav_{\init} \|^{-1} \) and \( (\fcl_{j}^{(0)}) \) by \( \| \betav_{\init} \| \) to ensure 
%\( \| \betav_{\init} \| = 1 \).
\end{mylist}

\paragraph{Step \( k \geq 0 \). The tensor \( \TEDR_{k} \) and the preliminary guess \( \betav_{\init} \)
are assumed available.}

\begin{mylist}


\item
For each \( j \in \JJJ \),
generate the sets of directions \( \Sph_{j} \) by drawing \( \sph \eqd \gaussv_{\TEDR_{k}}/\| \gaussv_{\TEDR_{k}} \| \)
for \( \gaussv_{\TEDR_{k}} \sim \ND(0,\TEDR_{k}^{2}) \);

\item 
For each \( j \in \JJJ \), \( \sph \in \Sph_{j} \), 
compute \( \Ima_{j,\sph}^{(k)} \, , S_{j,\sph}^{(k)} \, , \Uv_{j,\sph}^{(k)} \)
by \eqref{ytd7e7ydwuhd5eythdm} with
\begin{EQA}[c]
	\weight_{ij} 
	= 
	\weight_{ij}^{(k)}
	= 
	\Kern\bigl( \| \TEDR_{k} (\Xv_{i} - \xv_{j}) \|^{2} \bigr) 
	\, .
\end{EQA}

\item 
With 
\( \Imav_{j}^{(k)} = \bigl( \Ima_{j,\sph}^{(k)} \bigr)_{\sph \in \Sph_{j}} \),
\( \Sv_{j}^{(k)} = \bigl( S_{j,\sph}^{(k)} \bigr)_{\sph \in \Sph_{j}} \),
\( \UV_{j}^{(k)} = (\Uv_{j,\sph}^{(k)} \, , \sph \in \Sph_{j})^{\T} \),
solve
\begin{EQA}[c]
	\bigl\{ \betav_{k},(\fcn_{j}^{(k)},\fcl_{j}^{(k)})_{j \in \JJJ} \bigr\}
	=
	\arginf_{\betav \, , (\fcn_{j},\fclv_{j})_{j \in \JJJ}}
	\sum_{j \in \JJJ} \Bigl( \| \Imav_{j}^{(k)} - \fcn_{j} \Sv_{j}^{(k)} - \fcl_{j} \, \UV_{j}^{(k)} \betav \bigr\|^{2} \Bigr) 
%	+ \frac{\lambda}{2} \| \betav - \betav_{k-1} \|^{2}
	\, 
	\qquad
\end{EQA}
using alternating optimization from Section~\ref{SLLaltern} which starts at \( \betav_{\init} \).
%\item 
%%:
%Rescale \( \betav_{k} \) by \( \| \betav_{k} \|^{-1} \) and \( (\fcl_{j}^{(k)}) \) by \( \| \betav_{k} \| \) to ensure 
%\( \| \betav_{k} \| = 1 \).


\item
If \( \hiso_{k}/a < \hiso_{\min} \) then terminate with output \( \betav_{k} \).

\item
Otherwise

\begin{mylist}
\item
Update \( \betav_{\init} = \betav_{k} \).

\item 
Increase \( k \) by one.
Update \( \hiso_{k} = \hiso_{k-1} / a \).

\item 
Redefine \( \TEDR_{k} \) in the form \eqref{iufhvu8kjbmnll475edf} with \( \betav = \betav_{k-1} \) and \( \hiso = \hiso_{k} \):
\begin{EQA}[rcl]
	\TEDR_{k}^{2}
	&=&
	\hiso_{k}^{-2} \bigl( \aniso_{k}^{2} \, \Id_{\dimd} + \betav_{k-1} \betav_{k-1}^{\T} \bigr)
	\, ,
\label{iufhvu8kjbmnll475edfk}
\end{EQA}
where
\( \aniso_{k} \) is fixed by condition \eqref{jhgiverkytde4w34etycja}
with \( \TEDR = \TEDR_{k} \).
%\begin{EQA}[c]
%	\frac{1}{\nJJJ} \sum_{j \in \JJJ} \sumi \Kern\bigl( \| \TEDR_{k} (\Xv_{i} - \xv_{j}) \|^{2}
%	\bigr)
%	\geq 
%	\nmin
%	\, .
%	\qquad 
%\label{jhgiverkytde4w34etycja}
%\end{EQA}
\\
{\color{blue} plot \( \aniso_{k} \)}
\end{mylist}

\item
Loop.
% and \( (\fcn_{j}^{(k)},\fcl_{j}^{(k)})_{j \in \JJJ} \).
\end{mylist}


{\color{blue} 
\begin{mylist}
\item
Fix \( \dimd = 10 \) and explore the performance of the method 
for different values of tuning parameters, especially \( \nmin \), \( \nJJJ \), \( \nSph \), and \( \lambda \);

\item
behavior for different \( \sigma_{\eps} \), \( \sigmaX \);

\item
try with different \( \dimd \) and find a breakdown dimension \( \dimd \).

\end{mylist}
 
}



%%%%%%%%%%%%%%%%%%%%%%%%%%%%%%%  new section  %%%%%%%%%%%%%%%%%%%%%%%%%%%%%%%
\Section{Multi-index model}
\label{Smultind}


Consider the e.d.r. regression problem
\begin{EQA}
	Y_{i}
	&=&
	\gs(\Xv_{i}) + \eps_{i}
	=
	\fs(\PEDR \Xv_{i}) + \eps_{i}
	\, ,
\label{YifXieiedr}
\end{EQA}
where \( \Xv_{i} \in \R^{\dimd} \) for large \( \dimd \), 
\( \PEDR \) is an unknown element from the set \( \PROJ_{\dimd,\dimm} \) 
of orthogonal mappings from 
from \( \R^{\dimd} \) to \( \R^{\dimm} \) for \( \dimm \ll \dimd \), and
\( \fs \) is a function on \( \R^{\dimm} \).
The mapping \( \PEDR \) can be written as 
\begin{EQA}
	\PEDR \xv
	&=&
	(\betav_{1}^{\T} \xv) \betav_{1} + \ldots + (\betav_{\dimm}^{\T} \xv) \betav_{\dimm}
	\, ,
\label{ydgenb8jejohkou90tkr}
\end{EQA}
where \( \betav_{m}^{\T} \) are the orthonormal rows of \( \PEDR \).
Here the identifiability issue is even more severe than in the single-index model, because 
the set of the vectors \( \betav_{j} \) can be exchanged, rotated, etc. 

Similarly to the single-index case, 
the structural information about the model
can be recovered from the gradient vectors \( \nabla \gs(\xv) \).
More precisely, under \eqref{ucjcvjuvyvyf6r6rwjdy}, the \( \dimd \times \dimd \)-matrix
%by orthogonality of \( \betav_{1}, \ldots,\betav_{m} \)
\begin{EQA}[c]
	\JEDR
	\eqdef
	\sum_{j \in \JJJ} \nabla \gs(\xv_{j}) \, \nabla \gs(\xv_{j})^{\T}
	=
	\sum_{j \in \JJJ} \PEDR^{\T} \fclv_{j} \, \bigl( \PEDR^{\T} \fclv_{j} \bigr)^{\T}
%	\biggl( \sum_{m=1}^{\dimm} \fcl_{j,m} \, \betav_{m} \biggr)^{\T}
	=
	\PEDR^{\T} \biggl( \sum_{j \in \JJJ} \fclv_{j} \, \fclv_{j}^{\T} \biggr) \PEDR
	\, 
\label{ihjfiojihfer6t4rrt}
\end{EQA}
with \( \fclv_{j} = \nabla \fs(\PEDR \xv_{j}) \)
describes the multi-index structure.
Moreover, 
if all positive eigenvalues of \( \JEDR \) are different, 
then the SVD \( \JEDR = \PEDR^{\T} \Lambda \, \PEDR \) with \( \Lambda = \diag(\lambda_{1} \geq \ldots \geq \lambda_{\dimm}) \)
provides a unique canonical representation for the multi-index model \eqref{YifXieiedr}.


%%%%%%%%%%%%%%%%%%%% section %%%%%%%%%%%%%%%%%%%%
\Subsection{Preliminaries}
For a collection of centers \( \{ \xv_{j} \}_{j \in \JJJ} \),
and a localizing kernel \( \Kern \), consider the localizing weights of the form
\begin{EQA}
	\weight_{ij}
	&=&
	\Kern\bigl( \| \TEDR(\Xv_{i} - \xv_{j}) \|^{2} \bigr) 
	\, ,
\label{s7cuqkdwywdhbjwbfjwbdjm}
\end{EQA}
where \( \TEDR \) is a localizing tensor \( \TEDR \) (symmetric \( \dimd \)-matrix)
which reflects the available information about the multi-index structure in \eqref{YifXieiedr}.
At the beginning, an uninformative isotropic tensor \( \TEDR_{0} = \hiso_{0}^{-2} \Id_{\dimd} \) will be used.
During iteration, we learn the structure and use it to improve \( \TEDR \); see Section~\ref{SEDRsa} ahead.

For every \( j \in \JJJ \), assume a set \( \Sph_{j} \) or directions \( \sph \) (unit vectors in \( \R^{\dimd} \)) 
to be fixed.
Compute for all \( j \in \JJJ \), \( \sph \in \Sph_{j} \),
the quantities \( \Ima_{j,\sph}, S_{j,\sph} \), and \( \Uv_{j,\sph} \) 
by \eqref{ytd7e7ydwuhd5eythdm} as in the single-index case.
%\begin{EQ}[rcl]
%	\Ima_{j,\sph}
%	&=&
%	\sumi Y_{i} \langle \Xv_{i} - \xv_{j} \, ,\sph \rangle \, \weight_{ij} 
%	\, ,
%	\qquad
%	S_{j,\sph}
%	=
%	\sumi \langle \Xv_{i} - \xv_{j} \, ,\sph \rangle \, \weight_{ij}
%	\, ,
%%\label{ytd7e7ydwuhd5eythdYm}
%	\\
%	\Uv_{j,\sph}
%	&=&
%	\sumi (\Xv_{i} - \xv_{j}) \, \langle \Xv_{i} - \xv_{j} \, ,\sph \rangle \, \weight_{ij}
%	\in \R^{\dimd}
%	\, .
%	\qquad
%\label{ytd7e7ydwuhd5eythdm}
%\end{EQ}
Model equation \eqref{YifXieiedrSI} implies
\begin{EQA}
	\Ima_{j,\sph}
	&=&
	\sumi Y_{i} \langle \Xv_{i} - \xv_{j} \, ,\sph \rangle \, \weight_{ij}  
	\\
	&=&
	\sumi \fs(\PEDR \Xv_{i}) \langle \Xv_{i} - \xv_{j} \, ,\sph \rangle \, \weight_{ij}
	+ \sumi \eps_{i} \langle \Xv_{i} - \xv_{j} \, ,\sph \rangle \, \weight_{ij} \, .
\label{hgghjkklvddfrdshweym}
\end{EQA}
%
%as in \eqref{ytd7e7ydwuhd5eythdm}.
The function \( \fs \) is now \( \dimm \)-variate, and we consider a linear approximation 
\begin{EQA}
	\fs(\PEDR \xv) 
	& \approx &
	\fcn_{j} + \fclv_{j}^{\T} \PEDR (\xv - \xv_{j})
	=
	\fcn_{j} + \biggl( \sum_{m=1}^{\dimm} \fcl_{j,m} \, \betav_{m} \biggr)^{\T} (\xv - \xv_{j}) ,
\label{ucjcvjuvyvyf6r6rwjdy}
\end{EQA}
where \( \fclv_{j} = (\fcl_{j,m}) \in \R^{\dimm} \).
Correspondingly
\begin{EQA}[rcl]
	&& \nquad
	\sumi \fs(\PEDR \Xv_{i}) \, \langle \Xv_{i} - \xv_{j} \, ,\sph \rangle \, \weight_{ij} 
	\approx 
	\sumi \bigl\{ \fcn_{j} 
%	+ \sum_{m=1}^{\dimm} \fcl_{j,m} \, (\Xv_{i} - \xv_{j})^{\T} \betav_{m} \Bigr\}
	+ \fclv_{j}^{\T} \PEDR (\Xv_{i} - \xv_{j})
	\langle \Xv_{i} - \xv_{j} \, ,\sph \rangle \bigr\} \weight_{ij} 
	\\
	&=&
	\fcn_{j} S_{j,\sph} + \fclv_{j}^{\T} \PEDR \, \Uv_{j,\sph}
	=
	\fcn_{j} S_{j,\sph} + \Uv_{j,\sph}^{\T} \PEDR^{\T} \, \fclv_{j}
	\, .
\label{yd6x5ewbduwuqf64wfgfymi}
\end{EQA}
%where \( \fclv_{j} \cdot \betav = \fcl_{j,1} \, \betav_{1} + \ldots + \fcl_{j,\dimm} \, \betav_{\dimm} \).
The full dimensional parameter \( \prmtv \) includes
%the corresponding loading factors \( \mue_{1},\ldots,\mue_{\dimm} \), 
\( (\imav_{j,\sph}) \in \R^{\nJJJ \nSph} \), 
the coefficients \( (\fcn_{j},\fclv_{j})_{j \in \JJJ} \in \R^{\nJJJ(\dimm+1)} \),
and the target mapping \( \PEDR = (\betav_{1},\ldots,\betav_{\dimm})^{\T} \).
Denote \( \imav_{j} = \bigl( \ima_{j,\sph} \bigr)_{\sph \in \Sph_{j}} \in \R^{\nSph} \), 
\( \Imav_{j} = \bigl( \Ima_{j,\sph} \bigr)_{\sph \in \Sph_{j}} \in \R^{\nSph} \),
\( \Sv_{j} = \bigl( S_{j,\sph} \bigr)_{\sph \in \Sph_{j}} \in \R^{\nSph} \),
and let \( \UV_{j} \) be the \( \nSph \times \dimd \)-matrix with the rows \( \Uv_{j,\sph}^{\T} \),
\( \sph \in \Sph_{j} \).
Then the objective function reads
\begin{EQA}
	\LL(\prmtv)
	&=&
	\frac{1}{2} \sum_{j \in \JJJ} \Bigl( \| \Imav_{j} - \imav_{j} \|^{2} 
		+ \bigl\| \imav_{j} - \fcn_{j} \Sv_{j} - \UV_{j} \, \PEDR^{\T} \fclv_{j} \bigr\|^{2} \Bigr) 
	\, .
	\qquad
%\label{gtgyuy6ew5dtde32322jtmi}
\end{EQA}
The auxillary variables \( \zv_{j} \) can be eliminated, leading to
\begin{EQA}[rcl]
	\LL(\prmtv)
	&=&
	\LL(\PEDR,(\fcn_{j},\fclv_{j})_{j \in \JJJ})
	=
	\frac{1}{4} \sum_{j \in \JJJ} \bigl\| \Imav_{j} - \fcn_{j} \Sv_{j} - \UV_{j} \, \PEDR^{\T} \fclv_{j} \bigr\|^{2}	
	\, .
	\qquad
\label{gtgyuy6ew5dtde32322jtmiG0}	
\end{EQA}
This function should be minimized over the set of all coefficients \( (\fcn_{j},\fclv_{j})_{j \in \JJJ} \) and all
orthonormal mappings \( \PEDR \in \PROJ_{\dimd,\dimm} \) with \( \PEDR \PEDR^{\T} = \Id_{\dimm} \):
\begin{EQA}[c]
%	\prmtv_{0}
%	= 
%	\{ \PEDR_{0} \, , (\fcn_{j}^{(0)},\fclv_{j}^{(0)})_{j \in \JJJ} \} 
%	=
	\LL(\prmtv)
	\to
	\inf_{\PEDR \, , (\fcn_{j},\fclv_{j})_{j \in \JJJ}} 
	\, .
\label{ju7fudeju74334fhfd}
\end{EQA}
The class of mappings \( \PROJ_{\dimd,\dimm} \) is not convex; this makes the problem computationally hard.
We also face an identifiability issue: a solution \( \PEDR \) corresponds to 
a particular choice of the basis in the EDR space \( \EDR \).
This basis can be arbitrarily rotated.

%%%%%%%%%%%%%%%%%%%% section %%%%%%%%%%%%%%%%%%%%
\Subsection{Alternating optimization}
\label{SaltMI}
For solving \eqref{ju7fudeju74334fhfd}, one can apply the alternating procedure similar to Section~\ref{SLLaltern}.
The basic idea of the proposed approach is to relax the target \( \PEDR \in \PROJ_{\dimd,\dimm} \)
and optimize \( \LL(\prmtv) \) w.r.t. 
\( \BEDR = (\betav_{1},\ldots,\betav_{\dimm})^{\T} \in \R^{\dimm \times \dimd} \). 
Then the desired orthogonal mapping \( \PEDR \) will be obtained using SVD 
of the solution \( \BEDR \).
%
%To ensure numerical stability, we add a penalty term 
%\( \lambda \| \BEDR - \PEDR_{\init} \|_{\Fr}^{2} \),
%where \( \PEDR_{\init} \) is either a starting guess or 
%a previous step solution \( \PEDR_{\init} \).

The alternating optimization (AO) procedure assumes 
a starting guess \( \PEDR_{\init} \) to be given; see Section~\ref{SinitSI} ahead.
It repeats several times (\( k_{\max} \)) one-step alternating which reads as follows.

\paragraph{One-step alternating} 

\begin{mylist}
\item
With \( \PEDR_{\init} \) fixed, compute for every \( j \in \JJJ \)
\begin{EQA}[c]
	(\hat{\fcn}_{j},\hat{\fclv}_{j})
	=
	\arginf_{\fcn_{j},\fclv_{j}} \bigl\| \Imav_{j} - \fcn_{j} \Sv_{j} - \UV_{j} \PEDR_{\init}^{\T} \fclv_{j} \bigr\|^{2}
	\, .
\end{EQA}
{\color{blue} please provide a closed form; cf. Lemma~\ref{Lcjlj}.}

\item 
Then, with \( (\hat{\fcn}_{j},\hat{\fclv}_{j})_{j \in \JJJ} \) fixed, compute \( \hat{\BEDR} \):
\begin{EQA}[c]
	\hat{\BEDR}
	=
	\arginf_{\BEDR \in \R^{\dimm \times \dimd}} 
		\biggl\{ \sum_{j \in \JJJ} \bigl\| \Imav_{j} - \hat{\fcn}_{j} \Sv_{j} - \UV_{j} \BEDR^{\T} \hat{\fclv}_{j} \bigr\|^{2}
	+ \lambda \| \BEDR - \PEDR_{\init} \|_{\Fr}^{2} 
	\biggr\}	
	\, ,
\end{EQA}
{\color{blue} please provide a closed form; cf. Lemma~\ref{Lbeta}}

\item
Compute 
\begin{EQA}[rcl]
	\JEDR
	\eqdef
	\sum_{j \in \JJJ} (\hat{\BEDR}^{\T} \, \hat{\fclv}_{j}) \, \bigl( \hat{\BEDR}^{\T} \, \hat{\fclv}_{j} \bigr)^{\T}
	&=&
	\hat{\BEDR}^{\T} \Bigl( \sum_{j \in \JJJ} \hat{\fclv}_{j} \, \hat{\fclv}_{j}^{\T} \Bigr) \hat{\BEDR}
\end{EQA}
and its SVD
\begin{EQA}[c]
	\JEDR
	=
	\PEDR^{\T} \, \Lambda \, \PEDR
	\, ,
	\qquad
	\PEDR \, \PEDR^{\T}
	= 
	\Id_{\dimm}
	\, ,
	\qquad
	\Lambda = \diag(\lambda_{1} \geq \ldots \geq \lambda_{\dimm})
	\, .
\end{EQA}

\item 
Output \( \PEDR_{\outp} = \PEDR \) and \( \Lambda_{\outp} = \Lambda \).
\end{mylist}


\paragraph{Alternating optimization procedure} repeats
one-step alternating several times.
At every iteration, the output \( \PEDR_{\outp} \) of the previous step is used as input \( \PEDR_{\init} \).
The procedure delivers the last step results \( \PEDR_{\outp} \) and \( \Lambda_{\outp} \) after \( k_{\max} \) iterations.

{\color{blue} check one-step improvement and final quality of \( \PEDR \)}


%%%%%%%%%%%%%%%%%%%% section %%%%%%%%%%%%%%%%%%%%
\Subsection{Structural adaptation}
\label{SEDRsa}
The use of isotropic weights \( \weight_{ij} \) for computing the observables 
\( \Imav_{j} \), \( \Sv_{j} \), and \( \UV_{j} \) is suboptimal and leads to poor 
quality of recovering the EDR-space \( \EDR \).
The structure-adaptive approach from \cite{HJPS} suggests extracting the structural information from the current-step estimate 
and using it for improving the accuracy of the next-step estimate by applying 
anisotropic weights \( \weight_{ij} \).

Alternating procedure from Section~\ref{SaltMI} delivers 
a starting solution \( \PEDR_{0} \) of \eqref{ju7fudeju74334fhfd} and also the estimate \( \JEDR_{0} \) of \( \JEDR \)
from \eqref{ihjfiojihfer6t4rrt}.
These matrices can be used to improve the procedure by redefining the tensor \( \TEDR \) as
\begin{EQA}
	\TEDR^{2}
	&=&
	\hiso^{-2} \bigl( \aniso^{2} \, \Id_{\dimd} + \JEDR_{0} \bigr)
%	\sum_{m} \biggl( \sum_{j \in \JJJ} \fcl_{j,m}^{2} \biggr) \betav_{m} \betav_{m}^{\T} 
	\, ,
\label{yfijowjdo76dvuienk}
\end{EQA}
with \( \aniso < 1 \).
This value describes the degree of kernel anisotropy.
Alternatively, %one can use
\begin{EQA}[c]
	\TEDR^{2}
	=
	\hiso^{-2} \aniso^{2} (\Id_{\dimd} - \PEDR_{0}^{\T} \PEDR_{0}) + \hiso^{-2} \, \JEDR_{0}
	\, .
\end{EQA}
This tensor leads to anisotropic weights 
\( \weight_{ij} = \Kern\bigl( \| \TEDR(\Xv_{i} - \xv_{j}) \|^{2} \bigr) \)
which strongly penalize in the principal directions of the estimated EDR space.
It will also be used for renewal of the set of directions \( \Sph_{j} \) which are randomly sampled as
\( \gaussv_{\TEDR}/\| \gaussv_{\TEDR} \| \) for \( \gaussv_{\TEDR} \sim \ND(0,\TEDR^{2}) \). 
The sums \( \Ima_{j,\sph} \), \( S_{j,\sph} \), and \( \Uv_{j,\sph} \) in \eqref{ytd7e7ydwuhd5eythdm}
should be recomputed with these changes.
The improved estimator \( \tilde{\prmtv} \) is given by 
minimization of \( \LL(\prmtv) \) from \eqref{gtgyuy6ew5dtde32322jtmiG0}
after this update.
%
This structural adaptation step can be iterated several times.
The values \( \hiso, \aniso \) may change during iterations: 
\( \hiso \) decreases at every step with a fixed factor 
\( a = 2^{1/\dimm} \), while \( \aniso \) is fixed by condition similar to \eqref{jhgiverkytde4w34etycja}.

%%%%%%%%%%%%%%%%%%%% section %%%%%%%%%%%%%%%%%%%%
\Subsection{Initialization by local linear estimation}
\label{SinitSI}

Initialization is an important issue.
A natural proposal is to apply local linear gradient estimation
from Section~\ref{Siniloclin} and define 
\( \PEDR_{\init} \) using the first \( \dimm \) principal directions
for the set of the estimated gradients \( \hat{\nabla \gs}(\xv_{j}) \), \( j \in \JJJ \).
\\
{\color{blue} 
If \( \PEDRs \, \PEDR_{\init}^{\T} \approx 0 \) then no sense to continue.
Try with different \( \numb_{\lin} \).}

%%%%%%%%%%%%%%%%%%%% section %%%%%%%%%%%%%%%%%%%%
\Subsection{Implementation}
\label{SEDRmiproc}


Fix the set of tuning parameters; see Table~\ref{TdefMI}.
\begin{table}[h]
\begin{tabular}{c|c|c}
        \hline
        Parameter & Meaning & default \\
        \hline
        \( \nmin \) & \( \# \) of design points \( \Xv_{i} \) in a local vicinity & 10 \\
        \( \numb_{\lin} \) & parameter for local linear estimation & 
        between \( 2d \) and \( n/\nmin \) \\
        \( \nJJJ \) & size of \( \JJJ \) 				& between \( n/\nmin \) and \( n \) \\
		\( \nSph \) & size of each \( \Sph_{j} \)		& \( \min\{ \nmin, \dimd \} \) \\
%\( a (= 2^{1/\dimm}) \), 
%		\( \lambda & (= \nmin \dimm) \),
		\( \lambda \) & Lagrange multiplier 		& between 1 and \( \dimm \) \\
		\( \Kern \) & the kernel & Epanechnikov: \( \Kern(t) = (1-t^{2})_{+} \) \\
		\( a \)	& decrease factor for \( \hiso_{k} \) & \( 2^{1/\dimm} \) \\
		\( \hiso_{\min} \) & the smallest \( \hiso_{k} \) & \( 10 \sigmaX/n \) \\
		\hline
    \end{tabular}
\caption{Tuning parameters for multi-index procedure}    
\label{TdefMI}
\end{table}

%\( \dimm (\leq 3) \), 
%\( \nmin (= 10) \), 
%\( \nJJJ (= n) \), 
%\( \nSph (= \min\{ \nmin , \dimd \}) \),
%\( a (= 2^{1/\dimm}) \), 
%\( \lambda (= \dimm) \),
%\( k_{\max} (= 5) \).

\paragraph{Initialization.}
\begin{mylist}
\item
Generate \( (\xv_{j} \, , j \in \JJJ) \).

\item
Fix \( \hiso_{0} \) as the smallest value ensuring \eqref{jhgiverkytde4w34etycj}.
Define \( \TEDR_{0} = \hiso_{0}^{-2} \Id_{\dimp} \).

\item Set \( k=0 \).

\item
Fix \( \PEDR_{\init} \) by PCA on the set of local linear gradient estimates
\( \hat{\nabla \gs}(\xv_{j}) \).
%; see Section~\ref{Siniloclin}.
\end{mylist}


\paragraph{Step \( k \), the tensor \( \TEDR_{k} \) and the preliminary guess \( \PEDR_{\init} \) are assumed 
%\( \PEDR_{k-1} \) and \( \Lambda_{k-1} \) 
available.}

\begin{mylist}


\item
For each \( j \in \JJJ \), 
generate the sets of directions \( \sph \in \Sph_{j} \) with
\( \sph \eqd \gaussv_{\TEDR_{k}}/\| \gaussv_{\TEDR_{k}} \| \) for 
\( \gaussv_{\TEDR_{k}} \sim \ND(0,\TEDR_{k}^{2}) \).

\item 
Compute \( \bigl( \Ima_{j,\sph}^{(k)} \, , \, S_{j,\sph}^{(k)} \, , \, \Uv_{j,\sph}^{(k)} \bigr)_{j \in \JJJ , \, \sph \in \Sph_{j}} \) 
by \eqref{ytd7e7ydwuhd5eythdm} with \( \weight_{ij} = \Kern\bigl( \| \TEDR_{k} (\Xv_{i} - \xv_{j}) \|^{2} \bigr) \).

\item
With 
\( \Imav_{j}^{(k)} = \bigl( \Ima_{j,\sph}^{(k)} \bigr)_{\sph \in \Sph_{j}} \),
\( \Sv_{j}^{(k)} = \bigl( S_{j,\sph}^{(k)} \bigr)_{\sph \in \Sph_{j}} \),
\( \UV_{j}^{(k)} = (\Uv_{j,\sph}^{(k)} \, , \sph \in \Sph_{j})^{\T} \),
solve
\begin{EQA}[c]
	\bigl\{ \PEDR_{k},(\fcn_{j}^{(k)},\fclv_{j}^{(k)})_{j \in \JJJ} \bigr\}
	=
	\arginf_{\PEDR,(\fcn_{j},\fclv_{j})_{j \in \JJJ}}
	\sum_{j \in \JJJ} \bigl\| \Imav_{j}^{(k)} - \fcn_{j} \Sv_{j}^{(k)} - \UV_{j}^{(k)} \, \PEDR^{\T} \fclv_{j} \bigr\|^{2} 
%	+ \lambda \| \BEDR - \PEDR_{k-1} \|_{\Fr}^{2}
	\, 
\end{EQA}
%where \( \PEDR %= (\betav_{1}, \ldots, \betav_{\dimm})^{\T} \in \R^{\dimm \times \dimd} \), 
%\in \PROJ_{\dimd, \dimm} \), 
%\( \fcn_{j} \in \R \), \( \fclv_{j} \in \R^{\dimm} \), \( j \in \JJJ \),
using the alternating optimization as in Section~\ref{SaltMI} starting from 
\( \PEDR_{\init} \).

\item
Terminate if \( \hiso_{k}/a < \hiso_{\min} \).
Output \( \PEDR_{k} \), \( \Lambda_{k} \).
%, \( (\fcn_{j}^{(k)},\fclv_{j}^{(k)})_{j \in \JJJ} \).

\item 
%:
Otherwise:
\begin{mylist}


\item
Update \( \PEDR_{\init} = \PEDR_{k} \).

\item
Increase \( k \) by one and update \( \hiso_{k} = \hiso_{k-1} / a \).


\item
With \( \JEDR_{k-1} = \PEDR_{k-1}^{\T} \Lambda_{k-1} \PEDR_{k-1} \), define \( \TEDR_{k} \) by
\begin{EQA}[c]
	\TEDR_{k}^{2}
	\eqdef
	\aniso_{k}^{2} \, \hiso_{k}^{-2} \Id_{\dimd} + \hiso_{k}^{-2} \, \JEDR_{k-1}
	\quad
	\text{or }
	\quad
	\aniso_{k}^{2} \, \hiso_{k}^{-2} (\Id_{\dimd} - \PEDR_{k-1}^{\T} \PEDR_{k-1}) + \hiso_{k}^{-2} \, \JEDR_{k-1}
	\, .
\end{EQA}
For any \( \uv \in \R^{\dimd} \), this means
\begin{EQA}[rcl]
	\| \TEDR_{k} \uv \|^{2}
	&=&
	\aniso_{k}^{2} \, \hiso_{k}^{-2} \, \| \uv \|^{2} + \hiso_{k}^{-2} \, \| \Lambda_{k-1}^{1/2} \PEDR_{k-1} \uv \|^{2}
	\\
	\text{or }
	\quad
	\| \TEDR_{k} \uv \|^{2}
	&=&
	\aniso_{k}^{2} \, \hiso_{k}^{-2} \, \bigl( \| \uv \|^{2} - \| \PEDR_{k-1} \uv \|^{2} \bigr) 
	+ \hiso_{k}^{-2} \, \| \Lambda_{k-1}^{1/2} \PEDR_{k-1} \uv \|^{2}
	\, ,
\end{EQA}
{\color{blue} (please, implement and compare both versions; 
the second one is the default option)}
\\
where \( \aniso_{k} \) is the smallest value ensuring \eqref{jhgiverkytde4w34etycja}.
%\begin{EQA}[c]
%	\frac{1}{\nJJJ} \sum_{j \in \JJJ} \sumi 
%		\Kern\bigl( \| \TEDR_{k} (\Xv_{i} - \xv_{j}) \|^{2} \bigr)
%%	=
%%	\frac{1}{\nJJJ}\sum_{j \in \JJJ} \sumi \weight_{ij}
%	\geq 
%	\nmin
%	\, .
%\label{jhgiverkytde4w34etycjam}
%\end{EQA}
\end{mylist}
\item Loop.
\end{mylist}

{\color{blue} 
\begin{mylist}

\item
Parameters \( \nSph \), \( \lambda \), are especially questionable;

\item
check whether renewal of directions \( \sph \in \Sph_{j} \) is helpful.


\end{mylist}

}

%%%%%%%%%%%%%%%%%%%%%%%%%%%%%%%  new section  %%%%%%%%%%%%%%%%%%%%%%%%%%%%%%%
\Section{Manifold learning}

Consider a regression \( Y_{i} = \fs(\Xv_{i}) + \eps_{i} \) with a high-dimensional \( \Xv_{i} \in \R^{\dimd} \),
\( n = 1,\ldots,n \).
The manifold hypothesis assumes that the function \( \fs(\cdot) \) has a nearly multi-index structure in a local vicinity 
of each point \( \xv \in \R^{\dimd} \).
This means a local linear approximation
\begin{EQA}[c]
	\E Y_{i}
	\approx
	\fcn_{j} + \fclv_{j}^{\T} \PEDR_{j} (\Xv_{i} - \xv_{j})
\end{EQA}
for all points \( \Xv_{i} \) in a vicinity of \( \xv_{j} \).
Here not only the intercept \( \fcn_{j} \in \R \) and 
the slope coefficients \( \fclv_{j} = (\fcl_{1},\ldots,\fcl_{\dimm})^{\T} \in \R^{\dimm} \)
but also the mapping \( \PEDR_{j} = \PEDR(\xv) \in \PROJ_{\dimd,\dimm} \)
depend on \( \xv_{j} \).
%
The proposed approach extends the multi-index procedure from Section~\ref{Smultind} to this situation 
by considering a family \( (\PEDR_{j}) \) in place of one mapping \( \PEDR \) with a strong penalty on their variability.
This approach does not require any manifold description.
We also do not project the data on the manifold.

%%%%%%%%%%%%%%%%%%%% section %%%%%%%%%%%%%%%%%%%%
\Subsection{Multiscale modeling}
The manifold hypothesis means that for any fixed \( \xvd \) and any \( \xv \) close to \( \xvd \), 
one can use the mapping \( \PEDR(\xvd) \) in place of \( \PEDR(\xv) \). 
%
This discussion suggests introducing two scales.
One of them is for local linear approximation of the function \( \fs(\cdot) \) in the form
\begin{EQA}[c]
	\fs(\xv) - \fs(\xvd)
	\approx
	\fclv^{\T} \, \PEDR (\xv - \xvd) 
\end{EQA} 
with the coefficients \( \fclv = (\fcl_{m}) \) and mappings \( \PEDR \in \PROJ_{\dimd,\dimm} \) depending
on the central point \( \xvd \).
Another one for local manifold approximation:  
\(
	\PEDR(\xv)
	\approx
	\PEDR(\xvd)
%	\, .
\).
%A local tensor 
%\( \TEDR_{j} \) and bandwidth \( \hiso_{0} \) for the function \( \fs \),
%and tensor \( \TEDRm_{j} \) and bandwidth \( \hisom \) for the manifold.
At the initialization step of the procedure, we consider an isotropic vicinity
\( \{ \| \xv - \xvd \| \leq \hiso_{0} \} \) for approximation of \( \fs \)
and \( \{ \| \xv - \xvd \| \leq \hisom_{0} \} \) for continuity of \( \PEDR(\xv) \).
The approach assumes that the manifold is ``smoother'' than the function \( \fs \) and hence,
\( \hisom_{0} \gg \hiso_{0} \).

%, where \( \hiso_{0} \) is used for the weights \( \weight_{ij} \).
One can go one step further and assume that the manifold belongs to a linear subspace \( \EDR \) of a larger dimension \( \dimms \).
An example is a smooth curve in a two- or three-dimensional subspace.
Under such an assumption, one can first apply the multi-index procedure to recover this larger subspace \( \EDR \)
and then use this information for a warm start.

%%%%%%%%%%%%%%%%%%%% section %%%%%%%%%%%%%%%%%%%%
\Subsection{Objective function and manifold penalty}
The approach extends the procedure for the multi-index case from Section~\ref{SEDRmiproc}.
The basic idea behind the procedure from \eqref{ju7fudeju74334fhfd} is that the EDR space 
is the same for all \( \xv_{j} \).
This structural assumption does not hold in the manifold case.
However, it is nearly fulfilled if we restrict the analysis to a vicinity of each \( \xv_{j} \).
This suggests using two hierarchical localizing schemes described by 
localizing tensors \( \TEDR \) and \( \TEDRm \).

Suppose a collection of centers \( ( \xv_{j} )_{j \in \JJJ} \), and 
the sets of directions \( (\Sph_{j})_{j \in \JJJ} \) have been fixed.
% 
Further, let 
\( \Imav_{j} = \bigl( \Ima_{j,\sph} \bigr)_{\sph \in \Sph_{j}} \),
\( \Sv_{j} = \bigl( S_{j,\sph} \bigr)_{\sph \in \Sph_{j}} \),
\( \UV_{j}^{\T} = (\Uv_{j,\sph})_{\sph \in \Sph_{j}} \),
be computed as in \eqref{ytd7e7ydwuhd5eythdm} with 
\( \weight_{ij} = \Kern(\| \TEDR_{0} (\Xv_{i} - \xv_{j}) \|^{2}) \)
for \( \TEDR_{0} = \hiso_{0}^{-1} \Id_{\dimd} \) and a starting bandwidth \( \hiso_{0} \).
%
Now, for any \( \xv_{\jl} \in (\xv_{j}) \), consider a localized version of
\eqref{ju7fudeju74334fhfd}:
\begin{EQA}[rcl]
%	\prmtv^{(0)} 
	&& \nquad
	\bigl\{ \PEDR_{\jl} \, , (\fcn_{\jjl},\fclv_{\jjl})_{j \in \JJJ} \bigr\}
	= 
	\arginf_{\PEDR \, , (\fcn_{j},\fclv_{j})_{j \in \JJJ}} 
	\sum_{j \in \JJJ} 
	\bigl\| \Imav_{j} - \fcn_{j} \Sv_{j} - \UV_{j} \PEDR^{\T} \fclv_{j} \bigr\|^{2}
	\weightM_{\jjl}
%	\Kern\bigl( \| \TEDRm(\xv_{j} - \xv_{\jl}) \|^{2} \bigr)
%	+ \lambda_{\PEDR} \pent(\BEDR)
	\, ,
	\qquad \quad
\label{kjhviir8786tytcdeefm}
\end{EQA}
where \( \weightM_{\jjl} = \Kern(\| \TEDRm_{0} (\xv_{j} - \xv_{\jl}) \|^{2}) \) for
\( \TEDRm_{0} = \hisom_{0}^{-1} \Id_{\dimd} \) and \( \hisom_{0} \gg \hiso_{0} \).
The only difference between this problem and \eqref{ju7fudeju74334fhfd} is in using a multiplier (localizing weight)
\( \weightM_{\jjl} \)
%\( \Kern\bigl( \| \TEDRm(\xv_{j} - \xv_{\jl}) \|^{2} \bigr) \)
for each entry \( \bigl\| \Imav_{j} - \fcn_{j} \Sv_{j} - \UV_{j} \PEDR^{\T} \fclv_{j} \bigr\|^{2} \).
Such weighting restricts the multi-index approximation to a vicinity of \( \xv_{\jl} \) given by 
such weights \( (\weightM_{\jjl})_{j \in \JJJ} \). 
All local problems \eqref{kjhviir8786tytcdeefm} can be merged into one:
%leading to minimization of
\begin{EQA}[c]
	\sum_{\jl \in \JJJ} \sum_{j \in \JJJ}
	\bigl\| \Imav_{j} - \fcn_{\jjl} \Sv_{j} - \UV_{j} \PEDR_{\jl}^{\T} \fclv_{\jjl} \bigr\|^{2} \weightM_{\jjl}
	\to 
	\inf_{\{ \PEDR_{\jl} \, , (\fcn_{\jjl},\fclv_{\jjl})_{j \in \JJJ} \}_{\jl \in \JJJ}}
	\, .
	\qquad \quad
\label{kjhviir8786tytcdeefde}
\end{EQA}
%over \( \prmtv = \{ \PEDR_{\jl} \, , (\fcn_{\jjl},\fclv_{\jjl})_{j \in \JJJ} \}_{\jl \in \JJJ} \).
%
The choice of the family \( (\PEDR_{\jl}) \) can be stabilized and synchronized by a manifold penalty
\begin{EQA}[c]
	\pent_{\subM}(\PEDR)
	\eqdef
	\lambda \sum_{\jl \in \JJJ} \sum_{j \in \JJJ} \epsilon(\PEDR_{j} , \PEDR_{\jl})
%	\| \PEDR_{j} - \PEDR_{\jl} \|_{\Fr}^{2} 
	\; \weightM_{\jjl}
%	\Kern\bigl( \| \TEDRm_{j} (\xv_{j} - \xv_{\jl}) \|^{2} \bigr)
	\, .
\end{EQA}
The value \( \epsilon(\PEDR_{j} , \PEDR_{\jl}) \) measures discrepancy between \( \PEDR_{j} \) and \( \PEDR_{\jl} \).
The Frobenius or operator norms of the difference \( \PEDR_{j} - \PEDR_{\jl} \) cannot be used here because
each \( \PEDR_{j} \) is not uniquely defined; one can swap the sign of each row without changing the manifold projection. 
%The difference between projectors \( \PEDR_{j} \PEDR_{j}^{\T} - \PEDR_{\jl} \PEDR_{\jl}^{\T} \) is a reasonable alternative.
The basic idea behind the manifold penalty is that 
for close points \( \xv_{j} \) and  \( \xv_{\jl} \), the corresponding projectors  
\( \PEDR_{j}^{\T} \PEDR_{j} \) and \( \PEDR_{\jl}^{\T} \PEDR_{\jl} \) are also close to each other.  
This suggests
\begin{EQA}[c]
	\epsilon(\PEDR_{j} , \PEDR_{\jl})
	\eqdef
	\bigl\| \PEDR_{j} (\Id_{\dimd} - \PEDR_{\jl}^{\T} \PEDR_{\jl}) \bigr\|_{\Fr}^{2}
	\, .
\end{EQA}
This leads to minimization of %\( \LL_{\subM}(\prmtv) \) with
\begin{EQA}[c]
%	&& \nquad
%	(\BEDR_{\jl} \, , (\fcn_{\jjl},\fclv_{\jjl})_{j \in \JJJ} )_{\jl \in \JJJ}
%	\\
%	&=& 
%	\arginf_{(\BEDR_{j} \, , \fcn_{\jjl},\fclv_{\jjl})_{j \in \JJJ}} 
	\LL_{\subM}(\prmtv)
	\eqdef
	\sum_{\jl \in \JJJ} \sum_{j \in \JJJ}
	\bigl\| \Imav_{j} - \fcn_{\jjl} \Sv_{j} - \UV_{j} \PEDR_{\jl}^{\T} \fclv_{\jjl} \bigr\|^{2} \weightM_{\jjl}
	+ \lambda \sum_{\jl \in \JJJ} \sum_{j \in \JJJ} \epsilon( \PEDR_{j} , \PEDR_{\jl}) \; \weightM_{\jjl}
	\, .
\end{EQA}


%%%%%%%%%%%%%%%%%%%% section %%%%%%%%%%%%%%%%%%%%
\Subsection{Basic (one-step) procedure}
\label{SmaniEDR}
The global problem of minimizing the objective function \( \LL_{\subM}(\prmtv) \)
can be solved by alternating.
We assume that  
a preliminary guess \( \PEDR_{j,\init} \) is available for each point from the family \( (\xv_{j})_{j \in \JJJ} \).
Fix \( \jl \in \JJJ \) and consider a partial optimization problem:
\begin{EQA}[c]
	\sum_{j \in \JJJ}
	\bigl\| \Imav_{j} - \fcn_{\jjl} \Sv_{j} - \UV_{j} \PEDR_{\jl}^{\T} \fclv_{\jjl} \bigr\|^{2} \weightM_{\jjl}
	+ \lambda \sum_{j \in \JJJ} \epsilon( \PEDR_{\jl} , \PEDR_{j,\init}) \; \weightM_{\jjl}
	\to \inf_{\PEDR_{\jl} \, , (\fcn_{\jjl},\fclv_{\jjl})_{j \in \JJJ}}
	\, .
	\qquad
\label{iuhviiioijcwtyuvb}
\end{EQA}
It can be solved by the procedure from Section~\ref{SaltMI}
with obvious changes.
Given \( \PEDR_{\jl,\init} \), the coefficients \( (\fcn_{\jjl},\fclv_{\jjl})_{j \in \JJJ} \) can be obtained by soving 
the quadratic problem
\begin{EQA}[c]
	\bigl\| \Imav_{j} - \fcn_{\jjl} \Sv_{j} - \UV_{j} \PEDR_{\jl,\init}^{\T} \fclv_{\jjl} \bigr\|^{2}
	\to \inf_{(\fcn_{\jjl},\fclv_{\jjl})_{j \in \JJJ}}
	\, .
	\qquad
\label{kiiilfduy543rbggioe46}
\end{EQA}
Minimization of {\color{blue} \eqref{iuhviiioijcwtyuvb}} w.r.t. \( \PEDR_{\jl} \) is a bit more involved,
a relaxation of the constraint \( \PEDR_{\jl} \in \PROJ_{\dimd,\dimm} \) is helpful.
Define
\begin{EQA}[c]
	\IEDR_{\jl,\init}
	\eqdef
	\frac{1}{\WeightM_{\jl}} \sum_{j \in \JJJ} \PEDR_{j,\init}^{\T} \, \PEDR_{j,\init} \; \weightM_{\jjl}
	\, ,
	\qquad
	\WeightM_{\jl}
	\eqdef
	\sum_{j \in \JJJ} \weightM_{\jjl}	
	\, .
\label{jfu88uuujecvyyyyyue}
\end{EQA}
Note that \( \IEDR_{\jl,\init} \) is not a projector, because the set of projectors is not convex.
However, it is a subprojector in the sense that \( 0 \leq \IEDR_{\jl,\init} \leq \Id_{\dimd} \).



\begin{lemma}
\label{Lmanijjl}
With \( (\fcn_{\jjl}, \fclv_{\jjl})_{j \in \JJJ} \) fixed and \( \IEDR_{\jl,\init} \) from \eqref{jfu88uuujecvyyyyyue}, it holds
%the solution to \eqref{iuhviiioijcwtyuvb} coincides with the solution to
\begin{EQA}[rcl]
	&& \nquad
	\arginf_{\BEDR_{\jl} \in \R^{\dimm \times \dimd}} \biggl( 
		\sum_{j \in \JJJ} \bigl\| \Imav_{j} - \fcn_{\jjl} \Sv_{j} - \UV_{j} \BEDR_{\jl}^{\T} \fclv_{\jjl} \bigr\|^{2} \weightM_{\jjl}
		+ \lambda \sum_{j \in \JJJ} \bigl\| \BEDR_{\jl} (\Id_{\dimd} - \PEDR_{j,\init}^{\T} \PEDR_{j,\init}) \bigr\|_{\Fr}^{2} \; \weightM_{\jjl} 
	\biggr) 
	\\
	&=&
	\arginf_{\BEDR_{\jl} \in \R^{\dimm \times \dimd}} 
	\biggl( \sum_{j \in \JJJ}
		\bigl\| \Imav_{j} - \fcn_{\jjl} \Sv_{j} - \UV_{j} \BEDR_{\jl}^{\T} \fclv_{\jjl} \bigr\|^{2} \weightM_{\jjl}
		+ \lambda \WeightM_{\jl} \tr \bigl\{ \BEDR_{\jl}^{\T} \BEDR_{\jl} (\Id_{\dimd} - \IEDR_{\jl,\init} ) \bigr\}
	\biggr)
	\, .
	\qquad
\label{iuhviiioijcwtyuvbr}
\end{EQA}
\end{lemma}

\begin{proof}
{\color{blue} please check using the following lemma}
\end{proof}

\begin{lemma}
\label{LFrobML}
For any positive weights \( \weightM_{j} \) and projectors \( \PEDR_{j} \in \PROJ_{\dimd,\dimm} \) for \( j \in \JJJ \), it holds
\begin{EQA}[c]
	\sum_{j \in \JJJ} \| \BEDR (\Id_{\dimd} - \PEDR_{\jl}^{\T} \PEDR_{\jl}) \|_{\Fr}^{2} \, \weightM_{j}
	=
	\WeightM \tr\bigl\{ \BEDR^{\T} \BEDR (\Id_{\dimd} - \IEDR) \bigr\}
	\, ,
\end{EQA}
where
\begin{EQA}[c]
	\WeightM
	\eqdef
	\sum_{j \in \JJJ} \weightM_{j}
	\, ,
	\qquad
	\IEDR
	\eqdef
	\frac{1}{\WeightM} \sum_{j \in \JJJ} \PEDR_{j}^{\T} \PEDR_{j} \, \weightM_{j}
\end{EQA}
\end{lemma}

\begin{proof}
As \( \PEDR_{j}^{\T} \PEDR_{j} \) is a projector in \( \R^{\dimd} \) and\( \PEDR_{j} \, \PEDR_{j}^{\T} = \Id_{\dimm} \), 
 it holds 
\begin{EQA}[rcl]
	&& \nquad
	\sum_{j \in \JJJ} \| \BEDR (\Id_{\dimd} - \PEDR_{j}^{\T} \PEDR_{j}) \|_{\Fr}^{2}	 \, \weightM_{j}
	=
	\sum_{j \in \JJJ}\tr (\BEDR - \BEDR \, \PEDR_{j}^{\T} \PEDR_{j}) (\BEDR - \BEDR \, \PEDR_{j}^{\T} \PEDR_{j})^{\T} \weightM_{j}
	\\
	&=&
	\sum_{j \in \JJJ} \tr\Bigl\{ (\BEDR \, \BEDR^{\T} - \tr(\BEDR^{\T} \BEDR \, \PEDR_{j}^{\T} \PEDR_{j}) \Bigr\} \weightM_{j}
	=
	\WeightM \tr (\BEDR^{\T} \BEDR ) - \WeightM \tr (\BEDR^{\T} \BEDR \, \IEDR)
	\, .
\end{EQA}
This implies the assertion.
\end{proof}


Based on the solution \( \BEDR_{\jl} \) of \eqref{iuhviiioijcwtyuvbr}, compute
\begin{EQA}[rcl]
	\JEDR_{\jl}
	\eqdef
	\sum_{j \in \JJJ} (\BEDR_{\jl}^{\T} \, \fclv_{\jjl}) \, \bigl( \BEDR_{\jl}^{\T} \, \fclv_{\jjl} \bigr)^{\T}
		\weightM_{\jjl}
	&=&
	\BEDR_{\jl}^{\T} \biggl( 
		\sum_{j \in \JJJ} \fclv_{\jjl} \, \fclv_{\jjl}^{\T} \, \weightM_{\jjl} 
	\biggr) \BEDR_{\jl}
	\, .
\label{jhuifdyu8v8i6445hve}
\end{EQA}
As \( \BEDR_{\jl} \) is of rank \( \dimm \), the matrix \( \JEDR_{\jl} \) is 
also a rank \( \dimm \) matrix. 
Its SVD 
\begin{EQA}[c]
	\JEDR_{\jl}
	=
	\PEDR_{\jl}^{\T} \, \Lambda_{\jl} \, \PEDR_{\jl}
	\, ,
	\qquad
	\PEDR_{\jl} \, \PEDR_{\jl}^{\T}
	= 
	\Id_{\dimm}
	\, ,
	\qquad
	\Lambda_{\jl} = \diag(\lambda_{1} \geq \ldots \geq \lambda_{\dimm})
	\, 
	\qquad
\label{jhfuui887678ucfhie3}
\end{EQA}
delivers the mapping \( \PEDR_{\jl} \) which 
replaces \( \PEDR_{\jl,\init} \) and can be used for further steps
of the alternating procedure.
The matrices \( \JEDR_{\jl} \) describe the structure of the models and will be used 
for structural adaptation.

Now we can describe the \emph{one-step procedure}. 
Assume available:
%
\begin{mylist}
	\item 
	a collection of mappings \( (\PEDR_{\jl,\init})_{\jl \in \JJJ} \) from 
	initialization or the previous step;

	\item 
	a collection \( (\Imav_{j} \, , \Sv_{j} \, , \UV_{j})_{j \in \JJJ} \); see \eqref{ytd7e7ydwuhd5eythdm};

	\item 
%	manifold tensors \( \TEDRm_{\jl} \) 
%	and manifold weights \( \weightM_{\jjl} = \Kern(\| \TEDRm_{\jl} (\xv_{j} - \xv_{\jl}) \|^{2}) \); 
 	manifold weights \( \weightM_{\jjl} \).
\end{mylist}
%
\paragraph{One-step procedure.}
For each \( \jl \in \JJJ \):
%
\begin{mylist}
	\item
	Fix \( (\fcn_{\jjl},\fclv_{\jjl})_{j \in \JJJ} \) by solving \eqref{kiiilfduy543rbggioe46}.
%
	\item
	Compute \( \IEDR_{\jl,\init} \) by \eqref{jfu88uuujecvyyyyyue}.
%
	\item
	Solve
\begin{EQA}[c]
	\BEDR_{\jl}
	=
	\arginf_{\BEDR \in \R^{\dimm \times \dimd}} \biggl( \sum_{j \in \JJJ}
		\bigl\| \Imav_{j} - \fcn_{\jjl} \Sv_{j} - \UV_{j} \BEDR^{\T} \fclv_{\jjl} \bigr\|^{2} \weightM_{\jjl}
		+ \lambda_{\subM} \tr \bigl\{ \BEDR^{\T} \BEDR (\Id_{\dimd} - \IEDR_{\jl,\init}) \bigr\} 
	\biggr)
	\, .
	\qquad
\end{EQA}
%
	\item
	Compute \( \JEDR_{\jl} \) by \eqref{jhuifdyu8v8i6445hve} and \( \PEDR_{\jl} \) by SVD \eqref{jhfuui887678ucfhie3}.
	Update \( \PEDR_{\jl,\init} = \PEDR_{\jl} \).

\end{mylist}

%%%%%%%%%%%%%%%%%%%% section %%%%%%%%%%%%%%%%%%%%
\Subsection{Initialization by local linear estimation}
%The manifold learning procedure can be initialized using local linear gradient estimation.
Suppose a collection \( (\xv_{j})_{j \in \JJJ} \) to be fixed.
For each \( j \in \JJJ \), we can estimate the gradient \( \nabla \gs(\xv_{j}) \) 
by a local linear fit as suggested in Section~\ref{Siniloclin}.
With a proper choice of the parameter \( \nmin \), \( \numb_{\lin} \), and \( \nman \):
fix \( \hiso_{0} \), \( \hiso_{\lin} \), and \( \hisom_{0} \) as the smallest values ensuring 
\begin{EQ}[rcl]
	\frac{1}{\nJJJ}\sum_{j \in \JJJ} 
	\sumi \Kern\Bigl( \mfrac{\| \Xv_{i} - \xv_{j} \|^{2}}{\hiso_{0}^{2}} \Bigr)
	& \geq &
	\nmin
	\, ,
	\quad
	\frac{1}{\nJJJ} \sum_{j \in \JJJ} 
	\sumi \Kern\Bigl( \mfrac{\| \Xv_{i} - \xv_{j} \|^{2}}{\hiso_{\lin}^{2}} \Bigr)
	\geq 
	\numb_{\lin}
	\, ,
	\\
	\frac{1}{\nJJJ} \sum_{\jl \in \JJJ} 
	\sum_{j \in \JJJ} \Kern\Bigl( \mfrac{\| \xv_{j} - \xv_{\jl} \|^{2}}{\hisom_{0}^{2}} \Bigr)
	& \geq &
	\nman
	\, .
\label{hhfiecredgvhihiue86tML}
\end{EQ}
For each \( j \in \JJJ \), compute
\begin{EQA}[rcl]
	\begin{pmatrix}
		\hat{\gs}(\xv_{j}) \\
		\hat{\nabla \gs}(\xv_{j})
	\end{pmatrix}
	&=&
	\biggl\{ \sumi \binom{1}{\Xv_{i} - \xv_{j}} \, \binom{1}{\Xv_{i} - \xv_{j}}^{\T}
		\Kern\Bigl( \mfrac{\| \Xv_{i} - \xv_{j} \|^{2}}{\hiso_{\lin}^{2}} \Bigr)
	\biggr\}^{-1}
	\\
	&&
	\times
	\sumi Y_{i} \, \binom{1}{\Xv_{i} - \xv_{j}} \, 
		\Kern\Bigl( \mfrac{\| \Xv_{i} - \xv_{j} \|^{2}}{\hiso_{\lin}^{2}} \Bigr)
	\, .
\label{vgopo00d54323edfh}
\end{EQA}
Further, for each \( \jl \in \JJJ \), compute \( \JEDR_{\jl} \) as
\begin{EQA}[c]
	\JEDR_{\jl}
	\eqdef
	\sum_{j \in \JJJ} \hat{\nabla \gs}(\xv_{j}) \bigl\{ \hat{\nabla \gs}(\xv_{j}) \bigr\}^{\T} 
	\Kern\Bigl( \mfrac{\| \xv_{j} - \xv_{\jl} \|^{2}}{\hisom_{0}^{2}} \Bigr)
	\, ,
\label{jheew3regdw3bvio2ff}
\end{EQA}
and initialize \( \PEDR_{\jl,\init} \) by a PCA procedure using \( \dimm \) principal eigenvectors of \( \JEDR_{\jl} \).  
%Suppose that the points \( \xv_{\jl} \) are ordered with a possibly small distance between 
%any two subsequent points.
%The idea is to start with a non-informative guess for the first point \( \xv_{\jl} \) and solve \eqref{kjhviir8786tytcdeefm} 
%using the procedure from Section~\ref{SaltMI}. 
%The solution will be used as \( \PEDR_{\jl}^{(0)} \).
%Then we take the second point and solve \eqref{iuhviiioijcwtyuvb} using already obtained solutions \( \PEDR_{\jl}^{(0)} \)
%for manifold synchronization.
%The procedure should be continued until \( \PEDR_{\jl}^{(0)} \) is initialized for all \( (\xv_{\jl})_{\jl \in \JJJ} \).


%%%%%%%%%%%%%%%%%%%% section %%%%%%%%%%%%%%%%%%%%
\Subsection{Structure-adaptive manifold learning}
\label{SEDRmanproc}
This section describes the manifold learning procedure starting from initialization.



%\textbf{Tuning parameters.}
\begin{table}[b]
\begin{tabular}{c|l|c}
        \hline
        Parameter & Meaning & default \\
        \hline
        \( \dimm \) & manifold dimension & 1/2/3 \\
        \( \nmin \) & \( \# \) of design points \( \Xv_{i} \) in a local vicinity & 10 \\
		\( \nman \) & \( \# \) of the \( \xv_{j} \)'s in a manifold vicinity & \( \nJJJ/\nmin \) \\
        \( \nJJJ \) & size of \( \JJJ \) 				& between \( n/\nmin \) and \( n \) \\
		\( \nSph \) & size of each \( \Sph_{j} \)		& \( \min\{ \nmin, \dimd \} \) \\
%\( a (= 2^{1/\dimm}) \), 
%		\( \lambda & (= \nmin \dimm) \),
		\( \lambda_{\subM} \) & Lagrange multiplier 		& between 1 and \( \nmin \) \\
		\hline
    \end{tabular}
\caption{Tuning parameters for manifold learning procedure}    
\label{TdefML}
\end{table}

%\begin{mylist}
%\item
%\( \dimm (=1,2,3) \), 
%\item
%\( \nmin (= 10) \), 
%\item
%\( \nJJJ (= n/2) \), 
%\item
%\( \nman (= \nJJJ/10) \),
%\item
%\( \nSph (= \min\{ \nmin, \dimd \}) \),
%%\( a (= 2^{1/\dimm}) \), 
%\item
%\( \lambda (= \nmin \dimm) \),
%\item
%\( \lambda_{\subM} (=\nmin) \).
%\end{mylist}

\paragraph{Initialization.}
%Make preliminary calculations as follows:
\begin{mylist}

\item
Fix the set of tuning parameters; see Table~\ref{TdefML}.

\item
Generate the set \( (\xv_{j})_{j \in \JJJ} \).

\item
For each \( j \in \JJJ \), generate the set \( \Sph_{j} \) of unit vectors \( \sph \) in \( \R^{\dimd} \)
with \( |\Sph_{j}| = \nSph \).

\item
Fix \( \hiso_{0} \), \( \hiso_{\lin} \), and \( \hisom_{0} \) by \eqref{hhfiecredgvhihiue86tML}.
%as the smallest values ensuring
%\begin{EQA}[rcl]
%	\frac{1}{\nJJJ} \sum_{j \in \JJJ} 
%	\sumi \Kern\bigl( \hiso_{0}^{-2} \| \Xv_{i} - \xv_{j} \|^{2} \bigr)
%	& \geq &
%	\nmin
%	\, .
%\end{EQA}

\item 
For each \( j \in \JJJ \) and \( \sph \in \Sph_{j} \), compute  
\( \Ima_{j,\sph}^{(0)} \, , \, S_{j,\sph}^{(0)} \, , \, \Uv_{j,\sph}^{(0)} \)
by \eqref{ytd7e7ydwuhd5eythdm} with 
\begin{EQA}[c]
	\weight_{ij} 
	= 
	\weight_{ij}^{(0)}
	= 
	\Kern\bigl( \hiso_{0}^{-2} \| \Xv_{i} - \xv_{j} \|^{2} \bigr) 
	\, .
\end{EQA}

\item
For each \( j,\jl \in \JJJ \), compute \( \hat{\nabla \gs}(\xv_{j}) \) by \eqref{vgopo00d54323edfh} and 
\( \JEDR_{\jl} \) by \eqref{jheew3regdw3bvio2ff}.

\item
Define \( \PEDR_{\jl} \) using the first \( \dimm \) principal directions of \( \JEDR_{\jl} \).

\item
Run the one-step procedure several times (up to \( k_{\max} \)) using the output \( \PEDR_{\jl} \)
of the previous step as \( \PEDR_{\jl,\init} \) for all \( \jl \in \JJJ \).

\item
After the last step, set \( \PEDR_{\jl}^{(0)} = \PEDR_{\jl} \) for all \( \jl \in \JJJ \).
\end{mylist}
%After initialization, it is useful to carry out the one-step procedure several times (up to \( k_{\max} \))
%to make the first-step results more stable.
%After each run, the output \( (\PEDR_{\jl,\init})_{\jl \in \JJJ} \) should be used as the previous step collection 
%\( (\PEDR_{\jl}^{(0)})_{\jl \in \JJJ} \) for the next run.


%\paragraph{Sequential initialization.}
%For a given ordering on \( \JJJ \),
%compute sequentially the initial mappings \( \PEDR_{\jl}^{(0)} \), one after another, as follows:
%\begin{mylist}
%
%\item 
%Take the first point \( \xv_{\jl} \):
%\begin{mylist}
%
%\item
%Initialize \( \PEDR_{\jl}^{(0)} \) randomly or by the first-step global multi-index procedure, 
%set \( \BEDR_{\jl}^{(0)} = \PEDR_{\jl}^{(0)} \).
%
%\item With \( \BEDR_{\jl}^{(0)} \) fixed, apply the localized multi-index procedure to solve \eqref{kjhviir8786tytcdeefm}:
%\begin{mylist}
%
%	\item
%	Fix \( (\fcn_{\jjl},\fclv_{\jjl})_{j \in \JJJ} \) by solving \eqref{kiiilfduy543rbggioe46}.
%
%	\item
%	Fix \( \BEDR_{\jl} \) as a solution of
%\begin{EQA}[c]
%	\nquad
%%	\nquad
%%	\BEDR_{\jl}
%%	=
%	\inf_{\BEDR \in \R^{\dimm \times \dimd}} \biggl( \sum_{j \in \JJJ}
%		\bigl\| \Imav_{j}^{(0)} - \fcn_{\jjl} \Sv_{j}^{(0)} - \UV_{j}^{(0)} \BEDR^{\T} \fclv_{\jjl} \bigr\|^{2} \;
%	\weightM_{\jjl}^{(0)}
%		+ \lambda_{\subM} \| \BEDR  - \BEDR_{\jl}^{(0)} \|_{\Fr}^{2} 
%	\biggr)
%	\, .
%	\qquad
%\label{jjhidkmnvuy5rrtvnv}
%\end{EQA}
%
%	\item
%	Compute \( \PEDR_{\jl} \) and \( \Lambda_{\jl} \) by SVD of \( \JEDR_{\jl} \) as in \eqref{jhuifdyu8v8i6445hve} and
%	\eqref{jhfuui887678ucfhie3}.
%
%	\item
%	Update \( \PEDR_{\jl}^{(0)} = \BEDR_{\jl}^{(0)} = \PEDR_{\jl} \) and \( \Lambda_{\jl}^{(0)} = \Lambda_{\jl} \).
%	
%	\item 
%	Loop \( k_{\max} \) times.
%	
%\end{mylist}
%	
%	
%%	\item
%%	Set  \( \PEDR_{\jl}^{(0)} \) as solution \( \PEDR_{\jl} \) from the last step.
%	
%	\item
%	Define the set \( \JJJd = \{ \jl \} \).
%
%\end{mylist}
%
%\item
%Take the next point \( \xv_{\jl} \) with \( \jl \not\in \JJJd \).
%
%\begin{mylist}
%	\item
%	Compute \( \BEDR_{\jl}^{(0)} \) by extrapolation
%\begin{EQA}[c]
%	\BEDR_{\jl}^{(0)}
%	\eqdef
%	\biggl( \sum_{j \in \JJJd} \weightM_{\jjl} \biggr)^{-1} \sum_{j \in \JJJd} \PEDR_{j}^{(0)} \; \weightM_{\jjl}
%%	\, ,
%%	\qquad
%%	\WeightM_{\jl}
%%	\eqdef
%%	\sum_{j \in \JJJd} \weightM_{\jjl}	
%	\, .
%\label{jfu88uuujecvyyyyyued}
%\end{EQA}
%
%	\item 
%	Fix \( \BEDR_{\jl} \) by solving \eqref{jjhidkmnvuy5rrtvnv}.
%%\begin{EQA}[c]
%%	\BEDR_{\jl}
%%	=
%%	\arginf_{\BEDR \in \R^{\dimm \times \dimd}} \biggl( \sum_{j \in \JJJ}
%%		\bigl\| \Imav_{j} - \fcn_{\jjl} \Sv_{j} - \UV_{j} \BEDR^{\T} \fclv_{\jjl} \bigr\|^{2} \weightM_{\jjl}
%%		+ \lambda_{\subM} \| \BEDR  - \BEDR_{\jl}^{(0)} \|_{\Fr}^{2} 
%%	\biggr)
%%	\, .
%%	\qquad
%%\end{EQA}
%
%	\item
%	Compute \( \PEDR_{\jl} \) and \( \Lambda_{\jl} \) by SVD of \( \JEDR_{\jl} \) as in \eqref{jhuifdyu8v8i6445hve} and
%	\eqref{jhfuui887678ucfhie3}.
%
%	\item 
%	Set \( \PEDR_{\jl}^{(0)} = \PEDR_{\jl} \) and \( \Lambda_{\jl}^{(0)} = \Lambda_{\jl} \).
%	
%	\item
%	Add \( \jl \) to \( \JJJd \) and loop until \( \JJJd = \JJJ \).
%%The procedure also delivers \( \{ \JEDR_{\jl} \}_{\jl \in \JJJ} \).
%
%\end{mylist}
%\end{mylist}


Now we describe the step \( k \) of the main procedure for \( k \geq 1 \).


\paragraph{Step \( k \). Assumes the last-step solution 
\( (\PEDR_{\jl}^{(k-1)} \, , \Lambda_{\jl}^{(k-1)} )_{\jl \in \JJJ} \) available.}

\begin{mylist}

\item
Update \( \hiso_{k} = \hiso_{k-1}/a \), \( \hisom_{k} = \hisom_{k-1}/a \).


\item
With \( \JEDR_{\jl}^{(k-1)} = {\PEDR_{\jl}^{(k-1)}}^{\T} \Lambda_{\jl}^{(k-1)} \, \PEDR_{\jl}^{(k-1)} \),
define \( \TEDR_{j}^{(k)} \) and \( \TEDRm_{\jl}^{(k)} \) by
\begin{EQA}[rcl]
	(\TEDR_{j}^{(k)})^{2}
	&=&
	\frac{\aniso_{k}^{2}}{\hiso_{k}^{2}} \bigl( \Id_{\dimd} - {\PEDR_{j}^{(k-1)}}^{\T} \PEDR_{j}^{(k-1)} \bigr) 
	+ \frac{1}{\hiso_{k}^{2}} \, \JEDR_{j}^{(k-1)}
	\, ,
	\qquad
	j \in \JJJ
	\, ,
	\\
	(\TEDRm_{\jl}^{(k)})^{2}
	&=&
	\frac{\anisom_{k}^{2}}{\hisom_{k}^{2}} \bigl( \Id_{\dimd} - {\PEDR_{\jl}^{(k-1)}}^{\T} \PEDR_{\jl}^{(k-1)} \bigr) 
	+ \frac{1}{\hisom_{k}^{2}} \, \JEDR_{\jl}^{(k-1)}
	\, ,
	\qquad
	\jl \in \JJJ
	\, ,
\end{EQA}
where the anisotropy indices \( \aniso_{k} \) and \( \anisom_{k} \) are the largest values ensuring
\begin{EQA}[rcl]
	\frac{1}{\nJJJ} \sum_{j \in \JJJ} \sumi \Kern\bigl( \| \TEDR_{j}^{(k)} (\Xv_{i} - \xv_{j}) \|^{2} \bigr)
	& \geq &
	\nmin
	\, ,
	\\
	\frac{1}{\nJJJ} \sum_{\jl \in \JJJ} \sum_{j \in \JJJ} \Kern\bigl( \| \TEDRm_{\jl}^{(k)} (\xv_{j} - \xv_{\jl}) \|^{2} \bigr)
	& \geq &
	\nman
	\, .
\label{jhgiverkytde4w34etycjamM}
\end{EQA}
Denote for \( i = 1,\ldots,n \) and \( j,\jl \in \JJJ \)
\begin{EQA}[rcl]
	\weight_{ij}^{(k)} 
	&=& 
	\Kern\bigl( \| \TEDR_{j}^{(k)} (\Xv_{i} - \xv_{j}) \|^{2} \bigr) 
	\, ,
	\qquad
	\weightM_{\jjl}^{(k)} \eqdef \Kern\bigl( \| \TEDRm_{\jl}^{(k)} (\xv_{j} - \xv_{\jl}) \|^{2} \bigr)
	\, .
\end{EQA}

\item
For each \( j \in \JJJ \), 
update the sets of directions \( \Sph_{j} \) by drawing \( \sph \eqd \gaussv_{\TEDR_{j}^{(k)}}/\| \gaussv_{\TEDR_{j}^{(k)}} \| \)
for \( \gaussv_{\TEDR_{j}^{(k)}} \sim \ND(0,{\TEDR_{j}^{(k)}}^{2}) \).

\item 
For each \( j \in \JJJ \, \sph \in \Sph_{j} \), recompute
\( \Imav_{j}^{(k)} \), \( \Sv_{j}^{(k)} \), and \( \UV_{j}^{(k)} \) by \eqref{ytd7e7ydwuhd5eythdm} 
with \( \weight_{ij} = \weight_{ij}^{(k)} \).
%the weights  
%With %\( \BEDR = (\BEDR_{j} \, , j \in \JJJ) \) for 
%\( \BEDR_{j} = (\betav_{j,1}, \ldots, \betav_{j,\dimm})^{\T} \in \R^{\dimm \times \dimd} \), 

	\item
	For each \( \jl \in \JJJ \):
\begin{mylist}
	\item
	Fix \( (\fcn_{\jjl}^{(k)},\fclv_{\jjl}^{(k)})_{j \in \JJJ} \) by solving 
\begin{EQA}[c]
	(\fcn_{\jjl}^{(k)},\fclv_{\jjl}^{(k)})
	\eqdef
	\arginf_{(\fcn_{\jjl},\fclv_{\jjl})}
	\bigl\| \Imav_{j}^{(k)} - \fcn_{\jjl} \Sv_{j}^{(k)} - \UV_{j}^{(k)} {\PEDR_{\jl}^{(k-1)}}^{\T} \fclv_{\jjl} \bigr\|^{2}
	\, .
	\qquad
\label{kiiilfduy543rbgckr7r}
\end{EQA}
%	\eqref{kiiilfduy543rbggioe46}
%	with \( \Imav_{j}^{(k)} \, , \Sv_{j}^{(k)} \, , \UV_{j}^{(k)} \) in place of 
%	\( \Imav_{j} \, , \Sv_{j} \, , \UV_{j} \).
%
	\item
	Compute \( \IEDR_{\jl}^{(k-1)} \) as
\begin{EQA}[c]
	\IEDR_{\jl}^{(k-1)}
	\eqdef
	\biggl( \sum_{j \in \JJJ} \weightM_{\jjl}^{(k)} \biggr)^{-1} 
	\sum_{j \in \JJJ} \bigl( \PEDR_{j}^{(k-1)} \bigr)^{\T} \PEDR_{j}^{(k-1)} \; \weightM_{\jjl}^{(k)}
%	\, ,
%	\qquad
%	\WeightM_{\jl}
%	\eqdef
%	\sum_{j \in \JJJ} \weightM_{\jjl}	
	\, .
\label{jfu88uuujpvue2rchuu}
\end{EQA}

%
	\item
	Solve
\begin{EQA}[c]
	\BEDR_{\jl}^{(k)}
	=
	\arginf_{\BEDR \in \R^{\dimm \times \dimd}} \biggl( \sum_{j \in \JJJ}
		\bigl\| \Imav_{j}^{(k)} - \fcn_{\jjl}^{(k)} \Sv_{j}^{(k)} - \UV_{j}^{(k)} \BEDR^{\T} \fclv_{\jjl}^{(k)} \bigr\|^{2} 
		\weightM_{\jjl}^{(k)}
		+ \lambda_{\subM} \tr\bigl\{ \BEDR^{\T} \BEDR (\Id_{\dimd}  - \IEDR_{\jl}^{(k-1)} ) \bigr\}
	\biggr)
	\, .
	\qquad
\end{EQA}
%
	\item
	Compute \( \JEDR_{\jl}^{(k)} \) by \eqref{jhuifdyu8v8i6445hve} and 
	\( \PEDR_{\jl}^{(k)} \, , \Lambda_{\jl}^{(k)} \) by \eqref{jhfuui887678ucfhie3}
	with \( \BEDR_{\jl} = \BEDR_{\jl}^{(k)} \), \( \fclv_{\jjl} = \fclv_{\jjl}^{(k)} \).

\end{mylist}

%This step also delivers \( \JEDR_{\jl}^{(k)} \); see \eqref{jhuifdyu8v8i6445hve}.

\item
If \( \hiso_{k}/a \geq \hiso_{\min} \), increase \( k \) by one and loop.

\item
Otherwise, terminate.
Output \( \{ \PEDR_{\jl}^{(k)} \, , \Lambda_{\jl}^{(k)} \}_{\jl \in \JJJ} \).
\end{mylist}

{\color{blue} please, evaluate
\begin{mylist}

	\item 
	memory requirement as function of \( n \), \( \dimd \), \( \dimm \), \( \nJJJ \), \( \nSph \);
	\item
	computational cost as function of \( n \), \( \dimd \), \( \dimm \), \( \nJJJ \), \( \nSph \);
	\item
	accuracy of estimating the gradients \( \nabla \gs(\xv_{\jl}) \);
	\item
	accuracy of estimation after initialization (after the first run and then after each repetition)
	and frequency of a failure;
	\item
	one-step improvement by structural adaptation;
	\item
	choice of \( \lambda_{\manifold} \);
	\item
	degradation with growth of \( \dimd \).
\end{mylist}

}


%%%%%%%%%%%%%%%%%%%%%%%%%%%%%%%  new section  %%%%%%%%%%%%%%%%%%%%%%%%%%%%%%%
\Section{Average derivative procedure}
\label{Saverder}

%One limitation of the proposed approach is its memory requirements: the \( (\dimd+1) \)-matrix \( \KSe_{j} \)
%has to be stored at every point \( \xv_{j} \).
%
This section provides a light modification of the method 
which is focused on estimation of the index subspace \( \EDR \).


%%%%%%%%%%%%%%%%%%%% section %%%%%%%%%%%%%%%%%%%%
\Subsection{Single-index. Preliminaries}
\label{Saverdersi}
The single-index procedure from Section~\ref{Ssingind} recovers 
as the single-index vector \( \betav \), as the coefficients \( \fcn_{j}, \fcl_{j} \)
describing the value and slope coefficients of the function \( \fs \) at the point \( \xv_{j} \).
The average derivative approach focuses on estimation of directional derivatives. 
We use the notations of Section~\ref{Ssingind}.

Assumes a collection \( (\xv_{j})_{j \in \JJJ} \) is fixed.
Given a tensor \( \TEDR \), let \( \weight_{ij} = \Kern\bigl( \| \TEDR(\Xv_{i} - \xv_{j}) \|^{2} \bigr) \)
for all \( i=1,\ldots,n \) and \( j \in \JJJ \).
Also define the local mean
\begin{EQA}[c]
	\Xbar_{j}
	=
	\biggl( \sumi \weight_{ij} \biggr)^{-1}
	\sumi \Xv_{i} \, \weight_{ij}
	\, ,
	\qquad
	j \in \JJJ
	\, .
\label{uyhvvy883hierwfewAVE}
\end{EQA}
%The main difference with the approach of Section~\ref{Ssingind} is the use of directional averaging. 
For each point \( \xv_{j} \), suppose a set \( \Sph_{j} \) of directions (unit vectors) \( \sph \in \R^{\dimd} \) 
to be given.
Compute for every \( \sph \in \Sph_{j} \)
\begin{EQ}[rcl]
	\Ima_{j,\sph}
	& \eqdef &
	\sumi Y_{i} \, \langle \Xv_{i} - \Xbar_{j} \, , \sph \rangle \, \weight_{ij}
	\, ,
	\\
	\Uv_{j,\sph}
	& \eqdef &
	\sumi (\Xv_{i} - \Xbar_{j}) \, \langle \Xv_{i} - \Xbar_{j} \, ,\sph \rangle  \, \weight_{ij}
	\in \R^{\dimd}
	\, .
	\qquad
\label{juu8hjfioee433erhAVE}
\end{EQ}
The model equation \eqref{YifXieiedrSI} yields
\begin{EQA}[c]
	\E \Ima_{j,\sph}
	=
	\sumi \fs(\betav^{\T} \Xv_{i}) \, \langle \Xv_{i} - \Xbar_{j} \, ,\sph \rangle \, \weight_{ij}
	\, .
\end{EQA}
Moreover, an approximation \( \fs(t) \approx \fcn_{j} + \fcl_{j} (t - t_{j}) \) implies by definition of \( \Xbar_{j} \)
\begin{EQA}[rcl]
	\E \Ima_{j,\sph}
	& \approx &
	\sumi \bigl\{ \fcn_{j} + \fcl_{j} \, \betav^{\T} (\Xv_{i} - \xv_{j}) \bigr\} \, \langle \Xv_{i} - \Xbar_{j} \, ,\sph \rangle \, \weight_{ij}
	\\
	&=&
	\sumi \fcl_{j} \, \betav^{\T} (\Xv_{i} - \Xbar_{j}) \, \langle \Xv_{i} - \Xbar_{j} \, ,\sph \rangle \, \weight_{ij}
%	=
%	\fcl_{j} \, \uv_{j}^{\T} \KSme_{j} \betav
	=
	\fcl_{j} \, \Uv_{j,\sph}^{\T} \, \betav
	\, .
\end{EQA}
%where
%\begin{EQA}[c]
%	\Uv_{j,\sph}
%	\eqdef
%	\sumi (\Xv_{i} - \Sv_{j}) \, \langle \Xv_{i} - \Sv_{j} \, ,\sph \rangle  \, \weight_{ij}
%	\, .
%	\qquad \quad
%\label{uyhvvy883hierwfewAVE}
%\end{EQA}
%Now we define
%\begin{EQA}
%	\LL_{\lambda}(\prmtv)
%	&=&
%	\frac{1}{2} \| \Imav - \imav \|^{2} 
%	+ \frac{1}{2} \sum_{j \in \JJJ} \Bigl\| \imav_{j} - \fcl_{j} \, \Uv_{j}^{\T} \betav \Bigr\|^{2} 
%	+ \frac{\lambda}{2} \| \betav - \betav_{\init} \|^{2} \, .
%\label{tjjc7c6f5fAVE}
%\end{EQA}
With \( \Imav_{j} = (\Ima_{j,\sph})_{\sph \in \Sph_{j}} \in \R^{\nSph} \) and a \( \nSph \times \dimd \)-matrix 
\( \UV_{j} = (\Uv_{j,\sph})^{\T} \), this leads to the minimization of
\begin{EQA}[c]
	\sum_{j \in \JJJ} \sum_{\sph \in \Sph_{j}} \bigl( \Ima_{j,\sph} - \fcl_{j} \, \Uv_{j,\sph}^{\T} \, \betav \bigr)^{2} 
%	+ \lambda \| \betav - \betav_{\init} \|^{2}
	=
	\sum_{j \in \JJJ} \bigl\| \Imav_{j} - \fcl_{j} \, \UV_{j} \, \betav \bigr\|^{2}
%	+ \lambda \| \betav - \betav_{\init} \|^{2}
	\qquad
%	\to
%	\inf_{\betav,(\fcl_{j})_{j \in \JJJ}}
%	\, ,
\label{ikdfjoiefjoieuAVE}
\end{EQA}
over all unit vectors \( \betav \in \R^{\dimd} \) and \( (\fcl_{j})_{j \in \JJJ} \in \R^{\nSph} \). 
%which can be solved by alternating minimization as for \eqref{ikdfjoiefjoieuvbjmr3e}.

%%%%%%%%%%%%%%%%%%%% section %%%%%%%%%%%%%%%%%%%%
\Subsection{Single-index. Alternating optimization}
\label{SaverdersiAO}

Now we sketch the alternating optimization (AO) procedure for minimization of \eqref{ikdfjoiefjoieuAVE}.
It assumes available the vectors
\( \Imav_{j} = (\Ima_{j,\sph})_{\sph \in \Sph_{j}} \in \R^{\nSph} \) and \( \nSph \times \dimd \)-matrices
\( \UV_{j} = (\Uv_{j,\sph})^{\T} \) for \( j \in \JJJ \).
Also, a starting guess \( \betav_{\init} \) is required.

One-step alternating reads as foollows:

\begin{mylist}
\item
For \( \betav_{\init} \) fixed and each \( j \in \JJJ \), the value \( \fcl_{j} \) is estimated by
\begin{EQA}[c]
	\hat{\fcl}_{j}
	=
	\frac{\langle \Imav_{j}, \UV_{j} \, \betav_{\init} \rangle}{\| \UV_{j} \, \betav_{\init} \|^{2}}
	\, .
\end{EQA}

\item
For given \( (\hat{\fcl}_{j})_{j \in \JJJ} \), the vector \( \betav \) can be estimated as
\begin{EQA}[c]
	\hat{\betav}
	=
	\biggl( \sum_{j \in \JJJ} \hat{\fcl}_{j}^{\; 2} \; \UV_{j}^{\T} \UV_{j} + \lambda \Id_{\dimp} \biggr)^{-1}
	\biggl( \sum_{j \in \JJJ} \hat{\fcl}_{j} \, \UV_{j}^{\T} \Imav_{j} + \lambda \betav_{\init} \biggr)
	\, .
\end{EQA}
{\color{blue} please check.}

\item
The output \( \betav_{\outp} \) is obtained by normalizing \( \hat{\betav} \):
\begin{EQA}[c]
	\betav_{\outp}
	=
	\frac{1}{\| \hat{\betav} \|} \hat{\betav}
	\, .
\end{EQA}

\end{mylist}

\paragraph{Alternating optimization} can be viewed as iterating the one-step procedure.
At every iteration, the output \( \betav_{\outp} \) of the previous step is used as input \( \betav_{\init} \).
The procedure delivers the last step unit vector \( \betav_{\outp} \) after \( k_{\max} \) iterations.




%%%%%%%%%%%%%%%%%%%% section %%%%%%%%%%%%%%%%%%%%
\Subsection{Single-index. Initialization by local linear estimation}
\label{Saverdersini}

We aim at developing a data-driven approach. 
Initialization of the procedure is an important step.
We follow the approach of \cite{HJPS} based on local linear estimation.
Given \( \numb_{\lin} \ll n \), fix \( \hiso_{\lin} \) as the smallest value ensuring 
\begin{EQA}[c]
	\frac{1}{\nJJJ}\sum_{j \in \JJJ} 
	\sumi \Kern\Bigl( \mfrac{\| \Xv_{i} - \xv_{j} \|^{2}}{\hiso_{\lin}^{2}} \Bigr)
	\geq 
	\numb_{\lin}
	\, .
\label{hhfiecredgvhihiue86tade}
\end{EQA}
%The default value of \( \numb_{\lin} \) is \( n/\nmin \).
The proposal requires \( \numb_{\lin} > \dimd + 1 \).
The resulting bandwidth \( \hiso_{\lin} \) will be used 
for local linear approximation of the full dimensional function \( \gs(\xv) \).
Now, for each \( j \in \JJJ \), compute
\begin{EQA}[c]
	\begin{pmatrix}
		\hat{\gs}(\xv_{j}) \\
		\hat{\nabla \gs}(\xv_{j})
	\end{pmatrix}
	=
	\biggl\{ \sumi \binom{1}{\Xv_{i} - \xv_{j}} \, \binom{1}{\Xv_{i} - \xv_{j}}^{\T}
		\Kern\Bigl( \mfrac{\| \Xv_{i} - \xv_{j} \|^{2}}{\hiso_{\lin}^{2}} \Bigr)
	\biggr\}^{-1}
	\sumi Y_{i} \, \binom{1}{\Xv_{i} - \xv_{j}} \, 
		\Kern\Bigl( \mfrac{\| \Xv_{i} - \xv_{j} \|^{2}}{\hiso_{\lin}^{2}} \Bigr)
	\, .
\end{EQA}
Define \( \betav_{\init} \) as the first principal direction
for the set of the estimated gradients \( \{ \hat{\nabla \gs}(\xv_{j}) \}_{j \in \JJJ} \).
For this, compute
\begin{EQA}[c]
	\JEDR
	\eqdef
	\sum_{j \in \JJJ} \hat{\nabla \gs}(\xv_{j}) \{ \hat{\nabla \gs}(\xv_{j}) \}^{\T}
	\, .
\end{EQA}
Then \( \betav_{\init} \) is the eigenvector of \( \JEDR \) corresponding to its largest eigenvalue
\( \lambda_{\max}(\JEDR) \).


%%%%%%%%%%%%%%%%%%%% section %%%%%%%%%%%%%%%%%%%%
\Subsection{Single-index. Structure-adaptive procedure}
\label{SaverdersiP}

This section presents the single-index procedure, which consists of initialization and main iterations.
Initialization includes a few steps.

\begin{mylist}
\item
Fix the set of tuning parameters for the single-index procedure; see Table~\ref{TdefMIsi}.
\begin{table}[b]
\begin{tabular}{c|c|c}
        \hline
        Parameter & Meaning & default \\
        \hline
        \( \nmin \) & \( \# \) of design points \( \Xv_{i} \) in a local vicinity & 10 \\
        \( \numb_{\lin} \) & parameter for local linear estimation & 
        between \( 2d \) and \( n/\nmin \) \\
        \( \nJJJ \) & size of \( \JJJ \) 				& between \( n/\nmin \) and \( n \) \\
		\( \nSph \) & size of each \( \Sph_{j} \)		& \( \min\{ \nmin, \dimd \} \) \\
%\( a (= 2^{1/\dimm}) \), 
%		\( \lambda & (= \nmin \dimm) \),
		\( \lambda \) & Lagrange multiplier 		& 1 \\
		\( \Kern \) & the kernel & Epanechnikov: \( \Kern(t) = (1-t^{2})_{+} \) \\
		\( a \)	& decrease factor for \( \hiso_{k} \) & \( \sqrt{2} \) \\
		\( \hiso_{\min} \) & the smallest \( \hiso_{k} \) & \( 10 \sigmaX/n \) \\
		\hline
    \end{tabular}
\caption{Tuning parameters for the single-index procedure}    
\label{TdefMIsi}
\end{table}


\item
Fix a collection of points \( \xv_{j} \in \R^{\dimd} \), \( j \in \JJJ \).

\item
Identify the starting value \( \betav_{\init} \) by
the local linear procedure from Section~\ref{Saverdersini}.

\item
Fix \( \hiso_{0} \) as the smallest value ensuring
\begin{EQA}[c]
	\frac{1}{\nJJJ} \sum_{j \in \JJJ} 
	\sumi \Kern\bigl( \hiso_{0}^{-2} \| \Xv_{i} - \xv_{j} \|^{2} \bigr)
%	=
%	\frac{1}{\nJJJ} \sum_{j \in \JJJ} \sumi \weight_{ij}
	\geq 
	\nmin
	\, ;
\end{EQA}
{\color{blue} please plot \( \hiso_{0} \).}
\item
Define \( \TEDR_{0} = \hiso_{0}^{-1} \Id_{\dimd} \).
\end{mylist}




%\begin{mylist}
%\item
%Generate directions \( \sph \in \Sph_{j} \) uniformly on the unit sphere in \( \R^{\dimd} \).
%
%\item 
%for every \( j \in \JJJ \) and \( \sph \in \Sph_{j} \), compute \( \Ima_{j,\sph} \) and \( \Uv_{j,\sph} \) from 
%\eqref{uyhvvy883hierwfewAVE} and \eqref{juu8hjfioee433erhAVE} 
%with 
%\begin{EQA}[c]
%	\weight_{ij} = \Kern\bigl( \hiso_{0}^{-2} \| \Xv_{i} - \xv_{j} \|^{2} \bigr) 
%	\, .
%\end{EQA}
%\item 
%Find \( \betav_{\init} \) by minimizing \eqref{ikdfjoiefjoieuAVE} using alternating minimization
%with a random initialization.
%%\item 
%%%:
%%rescale \( \betav_{\init} \) by \( \| \betav_{\init} \|^{-1} \) and \( (\fcl_{j}^{(0)}) \) by \( \| \betav_{\init} \| \) to ensure 
%%\( \| \betav_{\init} \| = 1 \).
%\end{mylist}
%
%{\color{blue} check whether \( |\cos(\betav_{\init},\betavs)| \) is significantly positive}

\paragraph{Step \( k \geq 0 \). Assumed \( \TEDR_{k} \) and \( \betav_{\init} \) available.}

\begin{mylist}

\item
Generate the sets of directions \( \Sph_{j} \) using 
\( \sph \sim \gaussv_{\TEDR_{k}}/\| \gaussv_{\TEDR_{k}} \| \) for \( \gaussv_{\TEDR_{k}} \sim \ND(0,\TEDR_{k}^{2}) \).

\item 
For every \( j \in \JJJ \) and \( \sph \in \Sph_{j} \), compute \( \Ima_{j,\sph}^{(k)} \) and \( \Uv_{j,\sph}^{(k)} \) by \eqref{uyhvvy883hierwfewAVE} and \eqref{juu8hjfioee433erhAVE} 
with 
\begin{EQA}[c]
	\weight_{ij} 
	= 
	\weight_{ij}^{(k)}
	= 
	\Kern\bigl( \| \TEDR_{k} (\Xv_{i} - \xv_{j}) \|^{2} \bigr) 
	\, .
\end{EQA}

\item 
Solve by alternating minimization with initialization \( \betav_{\init} \):
\begin{EQA}[c]
	\{ \betav_{k},(\fcl_{j}^{(k)})_{j \in \JJJ} \} 
	=
	\arginf_{\betav,(\fcl_{j})_{j \in \JJJ}} 
	\sum_{j \in \JJJ} \bigl\| \Imav_{j}^{(k)} - \fcl_{j} \, \UV_{j}^{(k)} \betav \bigr\|^{2}
%	+ \lambda \| \betav - \betav_{k-1} \|^{2}
	\, .
\end{EQA}

{\color{blue} check whether \( |\cos(\betav_{\init},\betavs)| \) grows to one with iterations}
\item
If \( \hiso_{k} / a < \hiso_{\min} \), terminate. Output \( \PEDR_{k} \).
Otherwise,
\begin{mylist}
\item 
Update \( \betav_{\init} = \betav_{k} \).

\item
Increase \( k \) by 1.

\item 
Set \( \hiso_{k} = \hiso_{k-1} / a \).

\item 
Define 
\begin{EQA}[rcl]
	\TEDR_{k}^{2}
	&=&
	\hiso_{k}^{-2} \aniso_{k}^{2} \, \Id_{\dimd} + \hiso_{k}^{-2} \betav_{\init} \, \betav_{\init}^{\T}
	\, 
\label{iufhvu8kjbmnll475eAVE}
\end{EQA}
with \( \aniso_{k} \) fixed by the condition
\begin{EQA}[c]
	\frac{1}{\nJJJ} \sum_{j \in \JJJ} \sumi 
		\Kern\bigl( \| \TEDR_{k} (\Xv_{i} - \xv_{j}) \|^{2} \bigr)
	\geq 
	\nmin
	\, .
\label{jhgiverkytde4w34etyAVE}
\end{EQA}
{\color{blue} plot \( \aniso_{k} \)}
\end{mylist}


\item Loop.

%
\end{mylist}


%%%%%%%%%%%%%%%%%%%% section %%%%%%%%%%%%%%%%%%%%
\Subsection{Multiple index procedure}
\label{Saverdermi}
An extension to a multiple index procedure is straightforward.
Consider the e.d.r. regression problem
\begin{EQA}
	Y_{i}
	&=&
	\fs(\PEDR \Xv_{i}) + \eps_{i}
\label{YifXieiedrAVE}
\end{EQA}
where \( \Xv_{i} \in \R^{\dimd} \) for large \( \dimd \), 
\( \PEDR \) is an unknown element from the set \( \PROJ_{\dimd,\dimm} \) 
of orthogonal mappings from 
from \( \R^{\dimd} \) to \( \R^{\dimm} \) for \( \dimm \ll \dimd \), and
\( \fs \) is a function on \( \R^{\dimm} \).
The mapping \( \PEDR \) can be written as 
\begin{EQA}
	\PEDR \xv
	&=&
	(\betav_{1}^{\T} \xv) \betav_{1} + \ldots + (\betav_{\dimm}^{\T} \xv) \betav_{\dimm}
	\, ,
\label{ydgenb8jejAVE}
\end{EQA}
where \( \betav_{m}^{\T} \) are the orthonormal rows of \( \PEDR \).
Given a tensor \( \TEDR \), let the weights \( \weight_{ij} \) be defined as in \eqref{s7cuqkdwywdhbjwbfjwbdj}
for all \( i=1,\ldots,n \) and \( j \in \JJJ \).
The sums \( \Imav_{j} \) from \eqref{juu8hjfioee433erhAVE} can be expanded as
\begin{EQA}[rcl]
	\E \Ima_{j,\sph}
	& \approx &
	\sumi \biggl\{ \fcn_{j} + \sum_{\mm=1}^{\dimm} \fcl_{j,\mm} \, \betav_{\mm}^{\T} (\Xv_{i} - \xv_{j}) \biggr\} \, \langle \Xv_{i} - \Xbar_{j} \, ,\sph \rangle \, \weight_{ij}
	\\
	&=&
	\sum_{\mm=1}^{\dimm} \sumi \fcl_{j,\mm} \, \betav_{\mm}^{\T} (\Xv_{i} - \Xbar_{j}) \, \langle \Xv_{i} - \Xbar_{j} \, ,\sph \rangle \, \weight_{ij}
%	=
%	\fcl_{j} \, \uv_{j}^{\T} \KSme_{j} \betav
	=
	\Uv_{j,\sph}^{\T} \sum_{\mm=1}^{\dimm} \fcl_{j,\mm} \, \betav_{\mm}
	=
	\Uv_{j,\sph}^{\T} \PEDR^{\T} \fclv_{j}
	\, .
\end{EQA}
This suggests identifying the parameters \( (\fclv_{j})_{j \in \JJJ} \) and \( \PEDR \) by fitting the products 
\( \Uv_{j,\sph}^{\T} \PEDR^{\T} \fclv_{j} \) to \( \Ima_{j,\sph} \).
Therefore, each step of the procedure can be reduced to minimization of
\begin{EQA}[c]
	\sum_{j \in \JJJ} \sum_{\sph \in \Sph_{j}} 
	\bigl( \Ima_{j,\sph} - \Uv_{j,\sph}^{\T} \PEDR^{\T} \fclv_{j} \bigr)^{2}
%	+ \lambda \| \BEDR - \PEDR_{0} \|_{\Fr}^{2}
	=
	\sum_{j \in \JJJ}  
	\bigl\| \Imav_{j} - \UV_{j} \, \PEDR^{\T} \fclv_{j} \bigr\|^{2}
%	+ \lambda \| \BEDR - \PEDR_{0} \|_{\Fr}^{2}
	\, .
\label{ikdfjoiefvuheh3i}
\end{EQA}

\Subsubsection{One-step alternating}
As in the single-index case, minimization of \eqref{ikdfjoiefvuheh3i} can be performed 
by alternating optimization.
Given a preliminary guess \( \PEDR_{\init} \), we add a stabilizing penalty 
\( \lambda \| \BEDR - \PEDR_{\init} \|_{\Fr}^{2} \) 
and relax the constraint \( \PEDR \in \PROJ_{\dimd,\dimm} \) by 
replacing \( \PEDR \) with \( \BEDR \in \R^{\dimm \times \dimd} \).
This leads to the one-step problem
%
\begin{EQA}[c]
%	\sum_{j \in \JJJ} \sum_{\sph \in \Sph_{j}} 
%	\bigl( \Ima_{j,\sph} - \Uv_{j,\sph}^{\T} \BEDR^{\T} \fclv_{j} \bigr)^{2}
%	+ \lambda \| \BEDR - \PEDR_{0} \|_{\Fr}^{2}
%	=
	\inf_{\BEDR \, , \, (\fclv_{j})_{j \in \JJJ} }
	\sum_{j \in \JJJ}  
	\bigl\| \Imav_{j} - \UV_{j} \, \BEDR^{\T} \fclv_{j} \bigr\|^{2}
	+ \lambda \| \BEDR - \PEDR_{\init} \|_{\Fr}^{2}
	\, 
\label{ikdfjoiefjoieuAVEmi}
\end{EQA}
over the set \( \BEDR \in \R^{\dimm \times \dimd} \), \( (\fclv_{j})_{j \in \JJJ} \in \R^{\dimm \nSph} \).
The procedure minimizes the objective function w.r.t. the auxiliary parameters 
\(  (\fclv_{j})_{j \in \JJJ}\), then update \( \BEDR \) and projects on the space of 
orthogonal mappings.

\begin{mylist}
\item
Given \( \PEDR_{\init} \), for each \( j \in \JJJ \), compute
\begin{EQA}[c]
	\hat{\fclv}_{j}
	=
	\bigl( \PEDR_{\init} \, \UV_{j}^{\T} \UV_{j} \, \PEDR_{\init}^{\T} \bigr)^{-1} \PEDR_{\init} \, \UV_{j}^{\T} \, \Imav_{j}
	\, .
\end{EQA}

\item
Then compute
\begin{EQA}[c]
	\hat{\BEDR}
	=
	\arginf_{\BEDR} \biggl\{ \sum_{j \in \JJJ}  
	\bigl\| \Imav_{j} - \UV_{j} \, \BEDR^{\T} \hat{\fclv}_{j} \bigr\|^{2}
	+ \lambda \| \BEDR - \PEDR_{\init} \|_{\Fr}^{2} \biggr\}
	\, .
\end{EQA}
{\color{blue} please, provide a closed form solution}.
\item 
Compute
\begin{EQA}[rcl]
	\JEDR
	\eqdef
	\sum_{j \in \JJJ} (\hat{\BEDR}^{\T} \, \hat{\fclv}_{j}) \, \bigl( \hat{\BEDR}^{\T} \, \hat{\fclv}_{j} \bigr)^{\T}
	&=&
	\hat{\BEDR}^{\T} \Bigl( \sum_{j \in \JJJ} \hat{\fclv}_{j} \, \hat{\fclv}_{j}^{\T} \Bigr) \hat{\BEDR}
	\, 
\end{EQA}
and its SVD
\begin{EQA}[c]
	\JEDR
	=
	\PEDR^{\T} \Lambda \, \PEDR
	\, ,
	\qquad
	\PEDR \, \PEDR^{\T}
	= 
	\Id_{\dimm}
	\, ,
	\quad
	\Lambda = \diag(\lambda_{1} \geq \ldots \geq \lambda_{\dimm})
	\, .
\end{EQA}

\item 
Output \( \PEDR_{\outp} = \PEDR \) and \( \Lambda_{\outp} = \Lambda \).
\end{mylist}

This one-step procedure can be repeated several times.
At each step, the previous step solution \( \PEDR_{\outp} \) is used as \( \PEDR_{\init} \)
for the next step. 

\Subsection{Structure-adaptive procedure}
This section extends the single-index procedure to the multi-index case. 
This extension is straightforward and require only minor changes.
Initialization of the procedure includes a few steps.

\begin{mylist}
\item
Fix the set of tuning parameters for the single-index procedure; see Table~\ref{TdefMImi}.
\begin{table}[b]
\begin{tabular}{c|c|c}
        \hline
        Parameter & Meaning & default \\
        \hline
        \( \nmin \) & \( \# \) of design points \( \Xv_{i} \) in a local vicinity & 10 \\
        \( \numb_{\lin} \) & parameter for local linear estimation & 
        between \( 2d \) and \( n/\nmin \) \\
        \( \nJJJ \) & size of \( \JJJ \) 				& between \( n/\nmin \) and \( n \) \\
		\( \nSph \) & size of each \( \Sph_{j} \)		& \( \min\{ \nmin, \dimd \} \) \\
%\( a (= 2^{1/\dimm}) \), 
%		\( \lambda & (= \nmin \dimm) \),
		\( \lambda \) & Lagrange multiplier 		& 1 \\
		\( \Kern \) & the kernel & Epanechnikov: \( \Kern(t) = (1-t^{2})_{+} \) \\
		\( a \)	& decrease factor for \( \hiso_{k} \) & \( 2^{1/\dimm} \) \\
		\( \hiso_{\min} \) & the smallest \( \hiso_{k} \) & \( 10 \sigmaX/n \) \\
		\hline
    \end{tabular}
\caption{Tuning parameters for the single-index procedure}    
\label{TdefMImi}
\end{table}


\item
Fix a collection of points \( \xv_{j} \in \R^{\dimd} \), \( j \in \JJJ \).

\item
Identify the starting guess \( \PEDR_{\init} \) by
the local linear procedure from Section~\ref{Saverdersini}:
the resulting \( \PEDR_{\init} \) is defined by PCA on the set of estimated gradienst \( \hat{\gs}(\xv_{j}) \).
In the single-index case, onle the first principal direction is taken, now we
take the first \( \dimm \) principal directions.

\item
Fix \( \hiso_{0} \) as the smallest value ensuring
\begin{EQA}[c]
	\frac{1}{\nJJJ} \sum_{j \in \JJJ} 
	\sumi \Kern\bigl( \hiso_{0}^{-2} \| \Xv_{i} - \xv_{j} \|^{2} \bigr)
%	=
%	\frac{1}{\nJJJ} \sum_{j \in \JJJ} \sumi \weight_{ij}
	\geq 
	\nmin
	\, .
\end{EQA}

\item
Define \( \TEDR_{0} = \hiso_{0}^{-1} \Id_{\dimd} \).
\end{mylist}

\paragraph{Step \( k \geq 0 \). Assumed \( \TEDR_{k} \) and \( \PEDR_{\init} \) available.}

\begin{mylist}

\item
Generate the sets of directions \( \Sph_{j} \) using 
\( \sph \sim \gaussv_{\TEDR_{k}}/\| \gaussv_{\TEDR_{k}} \| \) for \( \gaussv_{\TEDR_{k}} \sim \ND(0,\TEDR_{k}^{2}) \).

\item 
For every \( j \in \JJJ \) and \( \sph \in \Sph_{j} \), compute \( \Ima_{j,\sph}^{(k)} \) and \( \Uv_{j,\sph}^{(k)} \) by \eqref{uyhvvy883hierwfewAVE} and \eqref{juu8hjfioee433erhAVE} 
with 
\begin{EQA}[c]
	\weight_{ij} 
	= 
	\weight_{ij}^{(k)}
	= 
	\Kern\bigl( \| \TEDR_{k} (\Xv_{i} - \xv_{j}) \|^{2} \bigr) 
	\, .
\end{EQA}

\item 
Solve by alternating minimization with initialization \( \PEDR_{\init} \):
\begin{EQA}[c]
	\{ \PEDR_{k},(\fcl_{j}^{(k)})_{j \in \JJJ} \} 
	=
	\arginf_{\PEDR,(\fcl_{j})_{j \in \JJJ}} 
	\sum_{j \in \JJJ} \bigl\| \Imav_{j}^{(k)} - \fcl_{j} \, \UV_{j}^{(k)} \betav \bigr\|^{2}
%	+ \lambda \| \betav - \betav_{k-1} \|^{2}
	\, .
\end{EQA}

\item
If \( \hiso_{k} / a < \hiso_{\min} \), terminate. Output \( \PEDR_{k} \) and \( \Lambda_{k} \).
Otherwise,
\begin{mylist}
\item 
Update \( \PEDR_{\init} = \PEDR_{k} \).

\item
Increase \( k \) by 1.

\item 
Set \( \hiso_{k} = \hiso_{k-1} / a \).

\item 
%With \( \JEDR_{k-1} = \PEDR_{k-1}^{\T} \Lambda_{k-1} \PEDR_{k-1} \), 
Define \( \TEDR_{k} \) by 
\begin{EQA}[rcl]
	\TEDR_{k}^{2}
	&=&
	\hiso_{k}^{-2} \aniso_{k}^{2} (\Id_{\dimd} - \PEDR_{k-1}^{\T} \PEDR_{k-1}) 
	+ \hiso_{k}^{-2} \, \PEDR_{k-1}^{\T} \Lambda_{k-1} \PEDR_{k-1}
	\, ,
\end{EQA}
where \( \aniso_{k} \) is fixed by the condition
\begin{EQA}[c]
	\frac{1}{\nJJJ} \sum_{j \in \JJJ} \sumi 
		\Kern\bigl( \| \TEDR_{k} (\Xv_{i} - \xv_{j}) \|^{2} \bigr)
	\geq 
	\nmin
	\, .
\end{EQA}
{\color{blue} plot \( \aniso_{k} \)}
\end{mylist}


\item Loop.

%
\end{mylist}


